%!TeX root=MemoriaTFG.tex

\chapter{Resultados}
\label{cap:resultados}

Este capítulo presenta los resultados de la evaluación de modelos
fundacionales dermatológicos y de las contribuciones originales del
trabajo, que concretan sus objetivos formales
(sección~\ref{sec:objetivos}). La secuencia es acumulativa: se verifica
primero la independencia entre los conjuntos de evaluación y el
preentrenamiento (sección~\ref{sec:auditoria-resultados}); sobre esa
base se estudia la transferencia supervisada (sondeo lineal, ajuste
fino, eficiencia en etiquetas), las tareas de segmentación,
clasificación \emph{zero-shot} y razonamiento multimodal, y los
aspectos de equidad y seguridad diagnóstica; y, finalmente, las
contribuciones propias sobre DermapixelAI~1.0 y Dermapixel R0, la
interpretabilidad mediante \emph{Sparse Autoencoders} y una síntesis
comparativa. En conjunto, los experimentos permiten analizar el
impacto del preentrenamiento específico de dominio, las distintas
estrategias de adaptación y el papel de los datos en la transferencia
de modelos fundacionales a un entorno clínico dermatológico en
castellano. La figura~\ref{fig:cap4-estructura} resume esta
organización.

\begin{figure}[H]
\centering
\small
\renewcommand{\arraystretch}{1.1}
\begin{tabular}{c}
\fbox{\textbf{Validez:} auditoría de solapamiento con Derm1M (§\ref{sec:auditoria-resultados})} \\
$\downarrow$ \\
\fbox{\textbf{Transferencia:} sondeo lineal, ajuste fino y eficiencia en etiquetas (datasets externos)} \\
$\downarrow$ \\
\fbox{\textbf{Segmentación} promptable basada en SAM/SAM2 (§\ref{sec:seg-results})} \\
$\downarrow$ \\
\fbox{\textbf{Zero-shot y modelos de lenguaje multimodales}} \\
$\downarrow$ \\
\fbox{\textbf{Equidad por fototipo y seguridad diagnóstica}} \\
$\downarrow$ \\
\fbox{\textbf{DermapixelAI~1.0 y Dermapixel R0} (validación en castellano y adaptación, O2/O3)} \\
$\downarrow$ \\
\fbox{\textbf{Interpretabilidad (SAE) y síntesis comparativa}} \\
\end{tabular}
\caption{Organización experimental del capítulo: de la validación de
la independencia de los datos a las contribuciones propias en
castellano.}
\label{fig:cap4-estructura}
\end{figure}

\section{Validez experimental: auditoría de solapamiento con Derm1M}
\label{sec:auditoria-resultados}

La validez de toda comparación de transferencia y \emph{zero-shot}
de este capítulo depende de que los conjuntos de test no se solapen
con el dataset de preentrenamiento; por ello el capítulo abre
comprobándolo. La auditoría intersecta los hashes MD5 de los doce
datasets externos con los $415\,451$ hashes únicos de la fracción
accesible de Derm1M e identifica un único dataset con solapamiento
exacto: \textbf{Dermnet}, cuyas $18\,302$ imágenes están todas
contenidas en Derm1M ($100\%$). Los once datasets restantes no
arrojan ninguna coincidencia. En particular, HIBA y MSKCC ---cuyo
origen en el ISIC archive motivaba una sospecha de procedencia---
presentan $0$ coincidencias MD5, por lo que esa sospecha no se
traduce en solapamiento exacto. El resultado completo por dataset
figura en el Anexo~\ref{cap:anexoa}.

El solapamiento de Dermnet no invalida las cifras que se reportan en
la sección~\ref{sec:rep-panderm} para PanDerm, DINOv2,
ConvNeXt-Large, EfficientNetV2-Large y SigLIP-Large, cuyo
preentrenamiento no incluye Derm1M. Sí condiciona, en cambio, la
interpretación de las cifras zero-shot de la familia DermLIP sobre
Dermnet (sección~\ref{sec:zs-results}), ya que DermLIP se entrenó
sobre Derm1M y vio durante el preentrenamiento la totalidad de ese
conjunto.

\section{Reproducción de PanDerm: sondeo lineal sobre diez datasets}
\label{sec:rep-panderm}

Esta sección evalúa si las representaciones de un modelo fundacional
dermatológico, sin reajustar el encoder, clasifican lesiones en
datasets externos no vistos. El codificador de mayor capacidad
transfiere de forma robusta entre datasets, con ventaja creciente en
los escenarios de mayor heterogeneidad de adquisición y complejidad
taxonómica, como cuantifican los diez datasets externos a continuación.

La tabla~\ref{tab:lp-panderm} aplica este protocolo a PanDerm Base y
Large sobre los diez datasets externos contrastados, conforme a
sección~\ref{sec:lp}. Se reportan cinco métricas: \emph{accuracy}
(Acc), \emph{balanced accuracy} (BAcc), AUROC \emph{one-vs-rest}, F1
ponderada (W-F1) y coeficiente kappa de Cohen.

\begin{table}[H]
\centering
\footnotesize
\setlength{\tabcolsep}{4pt}
\renewcommand{\arraystretch}{0.95}
\caption{Sondeo lineal con PanDerm Base y Large sobre diez datasets
externos. El asterisco (\textsuperscript{$\ast$}) marca Dermnet, el
único dataset con solapamiento exacto frente a Derm1M ($100\%$ por
hash MD5; sección~\ref{sec:auditoria-resultados}).
Mejores valores por dataset y métrica en negrita.}
\label{tab:lp-panderm}
\begin{tabular}{llrrrrr}
\hline
Dataset       & Modelo & Acc            & BAcc           & AUROC          & W-F1           & Kappa \\
\hline
HAM10000      & Base   & $0{,}853$      & $0{,}381$      & $0{,}892$      & $0{,}829$      & $0{,}709$ \\
HAM10000      & Large  & $\mathbf{0{,}888}$ & $\mathbf{0{,}575}$ & $\mathbf{0{,}954}$ & $\mathbf{0{,}880}$ & $\mathbf{0{,}814}$ \\
BCN20000      & Base   & $0{,}659$      & $0{,}292$      & $0{,}855$      & $0{,}615$      & $0{,}447$ \\
BCN20000      & Large  & $\mathbf{0{,}702}$ & $\mathbf{0{,}382}$ & $\mathbf{0{,}903}$ & $\mathbf{0{,}676}$ & $\mathbf{0{,}527}$ \\
PAD-UFES-20   & Base   & $0{,}725$      & $0{,}510$      & $0{,}913$      & $0{,}692$      & $0{,}369$ \\
PAD-UFES-20   & Large  & $\mathbf{0{,}772}$ & $\mathbf{0{,}642}$ & $\mathbf{0{,}949}$ & $\mathbf{0{,}760}$ & $\mathbf{0{,}527}$ \\
Dermnet$^\ast$& Base   & $0{,}450$      & $0{,}381$      & $0{,}888$      & $0{,}432$      & $0{,}437$ \\
Dermnet$^\ast$& Large  & $\mathbf{0{,}550}$ & $\mathbf{0{,}495}$ & $\mathbf{0{,}931}$ & $\mathbf{0{,}540}$ & $\mathbf{0{,}555}$ \\
WSI patches   & Base   & $0{,}781$      & $0{,}640$      & $0{,}976$      & $0{,}774$      & $0{,}842$ \\
WSI patches   & Large  & $\mathbf{0{,}868}$ & $\mathbf{0{,}792}$ & $\mathbf{0{,}991}$ & $\mathbf{0{,}871}$ & $\mathbf{0{,}896}$ \\
DDI           & Base   & $\mathbf{0{,}847}$ & $0{,}714$      & $0{,}827$      & $0{,}836$      & $0{,}482$ \\
DDI           & Large  & $\mathbf{0{,}847}$ & $\mathbf{0{,}764}$ & $\mathbf{0{,}860}$ & $\mathbf{0{,}846}$ & $\mathbf{0{,}535}$ \\
Derm7pt clín. & Base   & $0{,}762$      & $0{,}675$      & $0{,}808$      & $0{,}736$      & $0{,}397$ \\
Derm7pt clín. & Large  & $\mathbf{0{,}798}$ & $\mathbf{0{,}732}$ & $\mathbf{0{,}869}$ & $\mathbf{0{,}784}$ & $\mathbf{0{,}506}$ \\
Derm7pt dermo & Base   & $0{,}818$      & $0{,}749$      & $0{,}841$      & $0{,}811$      & $0{,}528$ \\
Derm7pt dermo & Large  & $\mathbf{0{,}836}$ & $\mathbf{0{,}766}$ & $\mathbf{0{,}858}$ & $\mathbf{0{,}829}$ & $\mathbf{0{,}570}$ \\
HIBA          & Base   & $0{,}883$      & $0{,}595$      & $0{,}894$      & $0{,}853$      & $0{,}275$ \\
HIBA          & Large  & $\mathbf{0{,}904}$ & $\mathbf{0{,}701}$ & $\mathbf{0{,}942}$ & $\mathbf{0{,}892}$ & $\mathbf{0{,}494}$ \\
MSKCC         & Base   & $0{,}721$      & $\mathbf{0{,}654}$ & $0{,}746$      & $0{,}720$      & $0{,}309$ \\
MSKCC         & Large  & $\mathbf{0{,}746}$ & $0{,}644$      & $\mathbf{0{,}751}$ & $\mathbf{0{,}732}$ & $\mathbf{0{,}316}$ \\
\hline
\emph{Media}  & Base   & $0{,}750$      & $0{,}559$      & $0{,}865$      & $0{,}730$      & --- \\
\emph{Media}  & Large  & $\mathbf{0{,}791}$ & $\mathbf{0{,}649}$ & $\mathbf{0{,}901}$ & $\mathbf{0{,}781}$ & --- \\
\hline
\end{tabular}
\end{table}

Sobre los diez datasets evaluados (estimación de semilla única,
sección~\ref{sec:estadistica}), PanDerm Large supera a PanDerm Base en
cuatro de las cinco métricas, con mejoras medias de $+4{,}1$\,pp en
accuracy, $+9{,}0$\,pp en \emph{balanced accuracy} y $+3{,}6$\,pp en
AUROC (figura~\ref{fig:lp-per-dataset}); la única excepción es la BAcc
de MSKCC ($-1{,}0$\,pp).

Las mayores ganancias aparecen en Dermnet, WSI~patches y PAD-UFES-20
---los escenarios más heterogéneos del benchmark---, lo que sugiere
que el aumento de capacidad del encoder beneficia especialmente la
generalización compleja. Dermnet, no obstante, está íntegramente
contenido en Derm1M ($100\%$ por hash MD5;
sección~\ref{sec:auditoria-resultados}) y su cifra debe leerse con ese
\emph{caveat}.

\begin{figure}[H]
\centering
\includegraphics[width=0.92\textwidth]{fig_lp_per_dataset.png}
\caption{Sondeo lineal con PanDerm Base y Large sobre los diez
datasets externos contrastados. Diferencia absoluta en BAcc
($+9{,}0$\,pp de promedio) y AUROC ($+3{,}6$\,pp de
promedio) a favor de PanDerm Large.}
\label{fig:lp-per-dataset}
\end{figure}

En conjunto, estos resultados indican que las representaciones
aprendidas por la familia PanDerm transfieren de forma consistente más
allá del dataset de origen, a múltiples cohortes externas, modalidades
de imagen y condiciones de adquisición. La ventaja de PanDerm Large se
observa en las tres modalidades evaluadas ---clínica, dermatoscópica e
histopatológica---, lo que sugiere que el aumento de capacidad del
encoder produce representaciones transferibles más allá de una única
modalidad de adquisición. Las medias por columna agregan datasets de
taxonomía, modalidad y prevalencia heterogéneas, por lo que se
interpretan como brújula descriptiva de la tendencia y no como
marcador absoluto del mejor modelo
(sección~\ref{sec:agregacion-datasets}).

\section{Benchmark comparativo entre familias}
\label{sec:benchmark}

Para distinguir si la ventaja de PanDerm Large refleja el
preentrenamiento dermatológico o una mera comparación favorable
frente a su versión Base, se sitúa frente a familias con filosofías
de aprendizaje distintas: arquitecturas convolucionales supervisadas,
modelos auto-supervisados generalistas y sistemas visión-lenguaje a
gran escala. La tabla~\ref{tab:benchmark} compara cinco modelos
representativos sobre HAM10000, PAD-UFES-20 y DDI bajo el mismo
protocolo de sondeo lineal.

\begin{table}[H]
\centering
\footnotesize
\setlength{\tabcolsep}{4pt}
\renewcommand{\arraystretch}{0.95}
\caption{Benchmark comparativo bajo sondeo lineal sobre HAM10000,
PAD-UFES-20 y DDI. Mejor valor por columna en negrita.}
\label{tab:benchmark}
\adjustbox{max width=\textwidth}{%
\begin{tabular}{lrrrrrrr}
\hline
Modelo                  & Tipo          & Params & HAM Acc       & HAM AUROC     & PAD Acc       & PAD AUROC     & DDI AUROC \\
\hline
ConvNeXt-Large          & ConvNet sup.  & $197$\,M & $0{,}883$ & $\mathbf{0{,}967}$ & $0{,}668$ & $0{,}894$ & $0{,}774$ \\
EfficientNetV2-Large    & ConvNet sup.  & $118$\,M & $0{,}860$ & $0{,}950$ & $0{,}690$ & $0{,}895$ & $0{,}766$ \\
DINOv2 ViT-L/14         & SSL general   & $307$\,M & $0{,}744$ & $0{,}847$ & ---       & ---       & --- \\
PanDerm Large           & SSL derm.     & $307$\,M & $0{,}888$ & $0{,}954$ & $\mathbf{0{,}772}$ & $\mathbf{0{,}949}$ & $\mathbf{0{,}860}$ \\
SigLIP-Large SO400M     & Vis-leng.\ gen.& $400$\,M & $\mathbf{0{,}900}$ & $0{,}971$ & $0{,}718$ & $0{,}903$ & $0{,}793$ \\
\hline
\end{tabular}}
\end{table}

La tabla~\ref{tab:gap-cnn} cuantifica la ventaja de PanDerm Large
sobre la mejor CNN supervisada en cada dataset.

\begin{table}[H]
\centering
\footnotesize
\setlength{\tabcolsep}{4pt}
\caption{PanDerm Large frente a la mejor CNN supervisada (la mejor de
ConvNeXt-L y EfficientNetV2-L para cada métrica).}
\label{tab:gap-cnn}
\begin{tabular}{lrrrrrr}
\hline
Dataset     & CNN Acc       & PanDerm Acc   & $\Delta$Acc (pp) & CNN AUROC     & PanDerm AUROC & $\Delta$AUROC (pp) \\
\hline
HAM10000    & $0{,}883$     & $0{,}888$     & $+0{,}5$         & $0{,}967$     & $0{,}954$     & $-1{,}3$ \\
PAD-UFES-20 & $0{,}690$     & $0{,}772$     & $+8{,}2$         & $0{,}895$     & $0{,}949$     & $+5{,}4$ \\
DDI         & $0{,}818$     & $0{,}847$     & $+2{,}9$         & $0{,}774$     & $0{,}860$     & $+8{,}6$ \\
\hline
\end{tabular}
\end{table}

PanDerm Large supera en accuracy a la mejor CNN supervisada sobre los
tres datasets (figura~\ref{fig:lp-benchmark}), con ventaja no uniforme:
$+0{,}5$\,pp en HAM10000 (dermatoscopia estandarizada) frente a
$+8{,}2$\,pp en PAD-UFES-20 (fotografía móvil) y $+2{,}9$\,pp en DDI.
HAM10000 constituye una excepción parcial: SigLIP-Large SO400M obtiene
la mejor accuracy y AUROC, y ConvNeXt-Large supera ligeramente a
PanDerm Large en AUROC ($-1{,}3$\,pp). Este comportamiento sugiere que
los beneficios del preentrenamiento de dominio son menores en entornos
dermatoscópicos altamente estandarizados; en cambio, PanDerm Large
encabeza en PAD-UFES-20 y DDI, los datasets clínicos más alejados del
dominio web.

\begin{figure}[H]
\centering
\includegraphics[width=0.88\textwidth]{fig_lp_comparison.png}
\caption{Benchmark comparativo bajo sondeo lineal sobre HAM10000,
PAD-UFES-20 y DDI. ConvNeXt-Large supera a PanDerm Large en AUROC
sobre HAM10000 ($-1{,}3$\,pp); PanDerm Large domina sobre
PAD-UFES-20 ($+8{,}2$\,pp Acc) y DDI ($+8{,}6$\,pp AUROC).}
\label{fig:lp-benchmark}
\end{figure}

Estos resultados muestran que no existe una jerarquía universal entre
familias de modelos: los sistemas visión-lenguaje generalistas pueden
igualar o superar a los especializados en escenarios dermatoscópicos
muy estandarizados, mientras que los modelos dermatológicos mantienen
una ventaja clara cuando aumenta la variabilidad clínica y de
adquisición. El preentrenamiento de dominio ofrece, por tanto, una
ventaja \emph{creciente} con la heterogeneidad del dominio, no una
superioridad absoluta ---lectura que anticipa la tesis de
especialización por tarea desarrollada en el
capítulo~\ref{cap:discusion}
(sección~\ref{sec:disc-especializacion}).

\section{Ajuste fino supervisado}
\label{sec:ft-results}

Esta sección delimita cuándo el ajuste fino del encoder mejora sobre
la representación congelada y cuándo es contraproducente, decisión
central para los datasets clínicos pequeños. El ajuste fino mejora con
una cantidad moderada o elevada de datos de entrenamiento y puede
degradar el rendimiento bajo baja disponibilidad de muestras. La
tabla~\ref{tab:ft-panderm} lo concreta comparando el
sondeo lineal con el ajuste fino supervisado de PanDerm Base sobre
cuatro datasets de tamaño dispar, con la configuración descrita
en sección~\ref{sec:ft}. Sobre HAM10000 se añade además
\emph{test-time augmentation} con cinco augmentaciones
deterministas.

\begin{table}[H]
\centering
\footnotesize
\setlength{\tabcolsep}{4pt}
\renewcommand{\arraystretch}{0.95}
\caption{Sondeo lineal (LP) frente a ajuste fino (FT) con PanDerm
Base. $\Delta$BAcc en puntos porcentuales respecto al LP.}
\label{tab:ft-panderm}
\begin{tabular}{llrrrrr}
\hline
Dataset      & Modo & Acc           & BAcc          & AUROC         & W-F1          & $\Delta$BAcc \\
\hline
HAM10000     & LP   & $0{,}853$     & $0{,}381$     & $0{,}892$     & $0{,}829$     & --- \\
HAM10000     & FT   & $\mathbf{0{,}920}$ & $\mathbf{0{,}852}$ & $\mathbf{0{,}978}$ & $\mathbf{0{,}923}$ & $+47{,}1$ \\
PAD-UFES-20  & LP   & $0{,}725$     & $0{,}510$     & $0{,}913$     & $0{,}692$     & --- \\
PAD-UFES-20  & FT   & $\mathbf{0{,}755}$ & $\mathbf{0{,}704}$ & $\mathbf{0{,}935}$ & $\mathbf{0{,}757}$ & $+19{,}4$ \\
DDI          & LP   & $\mathbf{0{,}847}$ & $\mathbf{0{,}714}$ & $\mathbf{0{,}827}$ & $\mathbf{0{,}836}$ & --- \\
DDI          & FT   & $0{,}774$     & $0{,}693$     & $0{,}682$     & $0{,}780$     & $-2{,}1$ \\
MSKCC        & LP   & $0{,}721$     & $0{,}654$     & $\mathbf{0{,}746}$ & $0{,}720$     & --- \\
MSKCC        & FT   & $\mathbf{0{,}731}$ & $\mathbf{0{,}684}$ & $0{,}719$     & $\mathbf{0{,}735}$ & $+3{,}0$ \\
\hline
\end{tabular}
\end{table}

El ajuste fino mejora sustancialmente el rendimiento en HAM10000
($+47{,}1$\,pp BAcc) y PAD-UFES-20 ($+19{,}4$\,pp), especialmente
sobre las clases minoritarias (melanoma, \emph{akiec}, \emph{df},
\emph{vasc}), responsables de la baja \emph{balanced accuracy} del
sondeo lineal. Sin embargo, el beneficio desaparece cuando disminuye
la disponibilidad de datos: DDI muestra una degradación global
compatible con sobreajuste y MSKCC sólo obtiene mejoras marginales. En
conjunto, los resultados sugieren que el ajuste fino resulta
ventajoso cuando existe una cantidad moderada o elevada de ejemplos
etiquetados, mientras que las representaciones congeladas son más
robustas en escenarios de datos limitados.

Este resultado justifica la exploración posterior de técnicas de
adaptación eficiente de parámetros (PEFT), que buscan capturar parte
de las ventajas del ajuste fino reduciendo el riesgo de sobreajuste en
datasets clínicos de tamaño moderado
(sección~\ref{sec:dermapixel-spanderm-v0}). Cabe asimismo cuantificar
cuántas etiquetas son realmente necesarias para aprovechar estas
representaciones, objeto de la sección siguiente.

\section{Eficiencia en etiquetas}
\label{sec:label-eff-results}

Se mide el rendimiento de PanDerm Large bajo submuestreos
estratificados crecientes del conjunto de entrenamiento, dado que el
etiquetado experto es el principal cuello de botella en dermatología.
Una fracción mínima del conjunto basta para casi saturar el
rendimiento. La tabla~\ref{tab:label-eff} lo cuantifica con
submuestreos al $1\%$,
$5\%$, $10\%$, $20\%$, $50\%$ y $100\%$ sobre HAM10000 y
PAD-UFES-20.

\begin{table}[H]
\centering
\footnotesize
\setlength{\tabcolsep}{4pt}
\renewcommand{\arraystretch}{0.95}
\caption{Eficiencia en etiquetas (sondeo lineal con PanDerm Large).
$N_{\mathrm{train}}$: tamaño efectivo del subconjunto de
entrenamiento. Mejor valor por columna en negrita.}
\label{tab:label-eff}
\begin{tabular}{lrrrrr}
\hline
$\%$ entrenam.  & $N_{\mathrm{train}}$ & Acc           & BAcc          & AUROC         & W-F1 \\
\hline
\multicolumn{6}{l}{\textit{HAM10000}} \\
$1\%$           & $82$                 & $0{,}850$     & $0{,}406$     & $0{,}918$     & $0{,}828$ \\
$5\%$           & $410$                & $0{,}870$     & $0{,}481$     & $0{,}932$     & $0{,}855$ \\
$10\%$          & $820$                & $0{,}875$     & $0{,}480$     & $0{,}939$     & $0{,}863$ \\
$20\%$          & $1\,641$             & $0{,}882$     & $0{,}547$     & $0{,}951$     & $0{,}873$ \\
$50\%$          & $4\,103$             & $\mathbf{0{,}894}$ & $\mathbf{0{,}589}$ & $0{,}952$ & $\mathbf{0{,}887}$ \\
$100\%$         & $8\,207$             & $0{,}888$     & $0{,}575$     & $\mathbf{0{,}954}$ & $0{,}880$ \\
\hline
\multicolumn{6}{l}{\textit{PAD-UFES-20}} \\
$1\%$           & $14$                 & $0{,}740$     & $0{,}599$     & $0{,}907$     & $0{,}727$ \\
$5\%$           & $74$                 & $0{,}751$     & $0{,}632$     & $0{,}919$     & $0{,}742$ \\
$10\%$          & $149$                & $0{,}738$     & $0{,}593$     & $0{,}923$     & $0{,}724$ \\
$20\%$          & $298$                & $0{,}748$     & $0{,}587$     & $0{,}931$     & $0{,}728$ \\
$50\%$          & $746$                & $0{,}757$     & $0{,}611$     & $0{,}941$     & $0{,}742$ \\
$100\%$         & $1\,493$             & $\mathbf{0{,}772}$ & $\mathbf{0{,}642}$ & $\mathbf{0{,}949}$ & $\mathbf{0{,}760}$ \\
\hline
\end{tabular}
\end{table}

La AUROC muestra una saturación temprana en ambos datasets. En
HAM10000, el $1\%$ de las imágenes de entrenamiento ($82$) ya recupera
una AUROC de $0{,}918$, mientras que el $5\%$ ($410$) alcanza el
$97{,}7\%$ del rendimiento máximo observado. Un comportamiento similar
aparece en PAD-UFES-20, donde $14$ imágenes permiten recuperar el
$95{,}5\%$ de la AUROC final. La \emph{balanced accuracy}, en cambio,
continúa mejorando con el aumento del conjunto de entrenamiento, lo que
refleja una mayor sensibilidad a las clases minoritarias.

Este comportamiento sugiere que gran parte del conocimiento visual
relevante ya ha sido adquirido durante el preentrenamiento, de modo
que el aprendizaje supervisado posterior actúa principalmente como un
mecanismo de calibración de la frontera de decisión: la mayor parte de
la capacidad discriminativa reside ya en el encoder preentrenado, lo
que reduce la dependencia de grandes volúmenes etiquetados en
escenarios clínicos. Las cifras proceden de una única partición y
semilla por nivel de submuestreo, por lo que delimitan una tendencia y
no un margen con cierre estadístico.

\section{Segmentación binaria de lesión}
\label{sec:seg-results}

Este bloque evalúa si un modelo fundacional de segmentación adaptado
al dominio dermatológico alcanza una delimitación de lesión de calidad
clínica y qué factores la determinan. La tabla~\ref{tab:seg} compara
SAM2.1-Large adaptado al dominio mediante \emph{fine-tuning} del
decodificador frente a SAM2 sin adaptar (\emph{zero-shot}), un
autoencoder convolucional sobre PanDerm Large (CAE-seg) y la
referencia publicada por Yan et~al.~\cite{panderm2025}, analizando
además su generalización fuera de dominio sobre ISIC2017 y PH2; las
subsecciones posteriores aíslan el efecto de la arquitectura, del
\emph{prompt} y de la variabilidad inter-anotador.

\begin{table}[H]
\centering
\footnotesize
\setlength{\tabcolsep}{4pt}
\renewcommand{\arraystretch}{0.95}
\caption{Segmentación binaria de lesión: Dice de test sobre tres
datasets dermatoscópicos. Modelos entrenados sobre ISIC2018 y
evaluados sin reentrenamiento sobre ISIC2017 y PH2.}
\label{tab:seg}
\begin{tabular}{lrrr}
\hline
Configuración                         & ISIC2018 Dice & ISIC2017 Dice & PH2 Dice \\
\hline
SAM2 \emph{zero-shot} (bbox de GT)    & ---           & $0{,}797$     & $0{,}886$ \\
CAE-seg (PanDerm-L) FT ISIC2018       & $0{,}894$     & $0{,}879$     & $0{,}924$ \\
PanDerm \cite{panderm2025} (ref.)     & $0{,}921$     & ---           & --- \\
\textbf{SAM2.1-L (decoder FT) ISIC2018}  & $\mathbf{0{,}947}$ & $\mathbf{0{,}945}$ & $\mathbf{0{,}960}$ \\
\hline
\end{tabular}
\end{table}

El detalle con la desviación estándar y el IoU sobre los datasets de
transferencia (ISIC2017 y PH2) figura en el Anexo~\ref{cap:anexob}
(tabla~\ref{tab:seg-detalle}).

SAM2.1-L con \emph{fine-tuning} del decodificador alcanza Dice
$0{,}947$ sobre ISIC2018, $+2{,}6$\,pp sobre la cifra publicada por
Yan et~al.~\cite{panderm2025} y $+5{,}3$\,pp sobre el CAE-seg: supera,
por tanto, la referencia previa. Su generalización fuera de dominio es
prácticamente nula (Dice $0{,}945$ en ISIC2017 y $0{,}960$ en PH2,
este último superior por su mayor limpieza visual). Y, sobre todo, la
adaptación del decodificador aporta $+14{,}8$\,pp en ISIC2017
($0{,}797 \to 0{,}945$) y $+7{,}4$\,pp en PH2 ($0{,}886 \to 0{,}960$)
frente a SAM2 sin adaptar (figura~\ref{fig:seg-generalization}): es la
adaptación al dominio, y no el tamaño bruto del modelo, lo que explica
la mayor parte de la mejora.

Desde una perspectiva aplicada, este resultado anticipa una lectura
que las subsecciones siguientes confirman: una vez disponible una caja
que \emph{localice} la lesión, la calidad del contorno depende poco
del segmentador concreto (sección~\ref{sec:seg-control-automatico}),
de modo que localizar la lesión puede importar más que la precisión
absoluta de la máscara.

\begin{figure}[H]
\centering
\includegraphics[width=0.88\textwidth]{fig_seg_generalization.png}
\caption{Segmentación binaria de lesión. SAM2.1-Large con
\emph{fine-tuning} del decodificador alcanza Dice $0{,}947$ sobre ISIC2018,
$+2{,}6$\,pp sobre el resultado publicado por
Yan~et~al.~\cite{panderm2025}, con transferencia OOD
prácticamente nula sobre ISIC2017 ($0{,}945$) y PH2
($0{,}960$).}
\label{fig:seg-generalization}
\end{figure}

\subsection{Comparativa de arquitecturas promptables}
\label{sec:seg-arquitecturas}

La cifra anterior mezcla generación de arquitectura, preentrenamiento
de dominio y tamaño de \emph{backbone}. Para separarlos, se evalúan de
forma pareada cinco arquitecturas \emph{promptables} sobre un dataset
dermatoscópico ampliado y deduplicado por hash MD5 ($14\,201$ imágenes
únicas de ISIC2018, ISIC2017, HAM10000-seg y PH2), bajo régimen común
(\emph{fine-tuning} denso del decodificador y el \emph{promptable
encoder}, encoder visual congelado, \emph{prompt} de \emph{bounding
box} de referencia; detalle de optimización en Cap.~3) y sobre el mismo
test canónico de ISIC2018 ($N = 1\,000$, blindado) que la cifra de
referencia $0{,}947$ de la tabla~\ref{tab:seg}. La
tabla~\ref{tab:seg-arquitecturas} reporta el Dice resultante.

\begin{table}[H]
\centering
\footnotesize
\setlength{\tabcolsep}{4pt}
\renewcommand{\arraystretch}{0.95}
\caption{Comparativa de cinco arquitecturas promptables sobre el
test canónico de ISIC2018 ($N = 1\,000$). Régimen común de
\emph{fine-tuning} del decodificador con \emph{prompt} de
\emph{bounding box}. Las diferencias de Dice son modestas en valor
absoluto pero significativas al test de Wilcoxon
pareado~\cite{wilcoxon1945}
($p < 10^{-6}$ en todos los contrastes).}
\label{tab:seg-arquitecturas}
\begin{tabular}{llllrr}
\hline
Modelo            & Arquitectura       & Preentr. & Params & Dice & Latencia \\
\hline
MedSAM2-tiny~\cite{medsam2} & SAM2~\cite{sam2}   & médico   & $39$\,M  & $\mathbf{0{,}9556}$ & $26$\,ms \\
SAM2.1-tiny       & SAM2~\cite{sam2}   & general  & $39$\,M  & $0{,}9517$ & $26$\,ms \\
SAM2.1-Large      & SAM2~\cite{sam2}   & general  & $224$\,M & $0{,}9503$ & $96$\,ms \\
SAM-Med2D~\cite{sammed2d} & SAM~\cite{sam2023} & médico   & $271$\,M & $0{,}9465$ & $20$\,ms \\
MedSAM v1~\cite{medsam}   & SAM~\cite{sam2023} & médico   & $94$\,M  & $0{,}9458$ & $74$\,ms \\
\hline
\end{tabular}
\end{table}

Tres conclusiones emergen de la comparación (todos los contrastes
pareados son significativos al test de Wilcoxon, $p < 10^{-6}$;
tabla~\ref{tab:seg-arquitecturas}). \textit{Primero}, la generación
SAM2 aporta más que el preentrenamiento médico: los tres modelos SAM2
superan a los basados en SAM v1, incluso a los de preentrenamiento
médico. \textit{Segundo}, el preentrenamiento médico ofrece una mejora
adicional, pero solo dentro de la misma generación y de menor
magnitud. \textit{Tercero}, el tamaño del \emph{backbone} aporta poco
en régimen \emph{box-prompted}: SAM2.1-tiny iguala o supera a
SAM2.1-Large con $5{,}7$ veces menos parámetros y menor latencia. La
mejor configuración es MedSAM2-tiny ($0{,}9556$, $26$\,ms), que
combina generación SAM2 y preentrenamiento médico en el
\emph{backbone} más compacto. Que la arquitectura pese más que el
preentrenamiento médico matiza la hipótesis central del trabajo: en
este régimen, el diseño del segmentador domina sobre la
especialización dermatológica. Las tres subsecciones siguientes acotan
cuánto del resultado depende del \emph{prompt}, su estabilidad y su
techo real.

\subsection{Valor del prompt: control automático sin caja}
\label{sec:seg-control-automatico}

Toda la comparativa anterior opera en régimen \emph{box-prompted}:
cada predicción parte de una \emph{bounding box} derivada de la
máscara de referencia. Para aislar la contribución de ese
\emph{prompt} se entrena un control automático que no recibe caja
alguna ---una U-Net con codificador ResNet-101 preentrenado en
ImageNet, ajuste fino denso de $51{,}5$\,M de parámetros sobre el
mismo dataset $v1c$--- y se evalúa sobre el test canónico de
ISIC2018. La tabla~\ref{tab:seg-automatico} contrasta su Dice con el
del mejor modelo \emph{promptable}.

\begin{table}[H]
\centering
\footnotesize
\setlength{\tabcolsep}{4pt}
\renewcommand{\arraystretch}{0.95}
\caption{Control automático \emph{frente a} régimen
\emph{promptable} sobre el test canónico de ISIC2018
($N = 1\,000$). La diferencia de Dice cuantifica el aporte
de la caja-\emph{prompt}.}
\label{tab:seg-automatico}
\begin{tabular}{llrrrr}
\hline
Modelo & Régimen & Dice & IoU & Precisión & Recall \\
\hline
MedSAM2-tiny      & \emph{promptable} & $\mathbf{0{,}9556}$ & $0{,}916$ & --- & --- \\
U-Net (ResNet-101) & automático        & $0{,}8811$ & $0{,}804$ & $0{,}856$ & $0{,}942$ \\
\hline
\end{tabular}
\end{table}

El control automático alcanza Dice $0{,}8811$, $7{,}5$\,pp por debajo
del mejor modelo \emph{box-prompted} ($0{,}9556$). La comparación de
magnitudes es reveladora: la arquitectura mueve $\sim 1{,}0$\,pp de
Dice (tabla~\ref{tab:seg-arquitecturas}) y el preentrenamiento médico
$\sim 0{,}4$\,pp, pero retirar la caja cuesta $7{,}5$\,pp ---entre
seis y siete veces más que cualquiera de los otros dos factores---. El
factor dominante de la segmentación dermatoscópica no es, por tanto,
ni la arquitectura ni el preentrenamiento, sino la disponibilidad de
una localización razonable de la lesión. El perfil de error del modo automático lo confirma
---recall alto ($0{,}942$) pero precisión desplomada ($0{,}856$) y
Hausdorff~95 entre $2{,}5$ y $3$ veces peor (desv.\ $0{,}121$), señal
de fallos catastróficos puntuales que la caja elimina---. El control
valida además que la comparativa fue justa: todos los modelos
compartían el mismo \emph{prompt}, de modo que sus diferencias
reflejan el segmentador y no la información de localización.

\subsection{Robustez del resultado: validación cruzada, prompt y calibración}
\label{sec:seg-robustez}

La cifra de $0{,}9556$ procede de una única partición. Una validación
cruzada de cinco pliegues sobre MedSAM2-tiny confirma su robustez:
Dice $0{,}9554 \pm 0{,}0010$ (IC95\% $[0{,}9543,\, 0{,}9564]$), con la
cifra principal dentro del intervalo y una desviación inter-pliegue
($8\times 10^{-4}$) que descarta dependencia de una partición
afortunada.

La estabilidad frente a la calidad del \emph{prompt} es crítica en
despliegue, donde la caja procede de un anotador o detector y no de la
máscara de referencia. La figura~\ref{fig:seg-box-robustness} reporta
el Dice bajo perturbaciones controladas. El \emph{jitter} de esquinas
($\pm 20$ px) apenas resta $0{,}86$\,pp, pero la escala es crítica y
asimétrica: apretar la caja es catastrófico ($-6{,}8$\,pp al $-10\%$,
$-15{,}4$\,pp al $-20\%$, por recortar lesión real) mientras que
aflojarla degrada poco ($-2{,}8$\,pp al $+10\%$, manteniendo Dice
$\approx 0{,}93$). La lectura operativa se desarrolla en
sección~\ref{sec:disc-cv-robustez}.

\begin{figure}[H]
\centering
\includegraphics[width=0.82\textwidth]{fig_seg_box_robustness.png}
\caption{Robustez de MedSAM2-tiny a la calidad de la caja-\emph{prompt}
sobre el test canónico de ISIC2018 ($N = 1\,000$). El \emph{jitter}
de esquinas es prácticamente inocuo ($-0{,}86$\,pp), pero la escala
es crítica y asimétrica: apretar la caja degrada mucho más que
aflojarla.}
\label{fig:seg-box-robustness}
\end{figure}

El resto de comprobaciones confirman un modelo bien comportado y no
alteran la conclusión: el ajuste de bajo rango (LoRA
$r\in\{4,8,16,32\}$) iguala al ajuste denso del decodificador; la
calibración a nivel de píxel es buena (ECE $0{,}0082$); la equidad por
fototipo (ITA) muestra una brecha de solo $1{,}1$\,pp entre fototipos
extremos; y el aumento en test y el \emph{ensemble} de pliegues aportan
mejoras marginales ($+0{,}13$/$+0{,}14$\,pp) que no compensan su coste,
por lo que se recomienda un único modelo.

\subsection{El techo de la métrica: acuerdo inter-anotador}
\label{sec:seg-techo-anotador}

El hallazgo más relevante de este bloque no es la cifra absoluta de
Dice, sino su relación con el límite humano. El Dice se mide contra
una máscara de referencia trazada por una persona, y el contorno de
una lesión es intrínsecamente ambiguo: sobre el subconjunto de IMA++
con anotación múltiple ($312$ imágenes con dos o más máscaras
independientes, reservado y nunca usado en entrenamiento), el acuerdo
\emph{inter-anotador} es de solo $0{,}728$, mientras que el modelo
concuerda con el consenso en $0{,}915$ ---es decir, coincide con el
consenso mejor de lo que los propios humanos coinciden entre sí--- y
cae dentro de la envolvente de las máscaras humanas el $97{,}8\%$ de
las veces.

La consecuencia es interpretativa y profunda: las diferencias de
$\sim 1$\,pp entre arquitecturas (tabla~\ref{tab:seg-arquitecturas})
son menores que el desacuerdo entre anotadores humanos, de modo que
perseguir mejoras de décimas de Dice tiene escaso sentido práctico ---
el resultado ya se sitúa en el techo que la propia tarea permite. El modelo es
algo sobreconfiado en bordes ambiguos: su error de calibración contra
la máscara de consenso blanda sube a $0{,}094$, frente al $0{,}0082$
contra una referencia binaria. Dentro del alcance evaluado
---segmentación binaria de lesión sobre imagen dermatoscópica, con
\emph{prompt} de caja y test blindado de ISIC2018--- la principal
limitación reside ya en la incertidumbre de la anotación humana, no en
el segmentador; el resultado no se extiende a la segmentación
multilesión, a la adquisición clínica no dermatoscópica ni al régimen
sin caja-\emph{prompt} (sección~\ref{sec:seg-control-automatico}). La
supervisión en forma de lenguaje, en cambio, motiva el bloque
siguiente sobre clasificación \emph{zero-shot} multimodal.

\section{Clasificación zero-shot multimodal}
\label{sec:zs-results}

Esta sección analiza de qué depende que un modelo visión-lenguaje
clasifique bien sin entrenar cabeza alguna, régimen que prometería
despliegue inmediato sin etiquetas. El factor crítico no es la imagen
sino la compatibilidad entre el tokenizador del texto y el espacio del
encoder. La tabla~\ref{tab:zs-ham} lo cuantifica sobre HAM10000
($7$ clases) con tres configuraciones de \emph{prompts} y dos
tokenizadores distintos.

\begin{table}[H]
\centering
\footnotesize
\setlength{\tabcolsep}{4pt}
\renewcommand{\arraystretch}{0.95}
\caption{Zero-shot multimodal sobre HAM10000 ($N_{\mathrm{test}} = 1\,232$,
$7$ clases). Mejor valor por columna en negrita.}
\label{tab:zs-ham}
\begin{tabular}{llllrr}
\hline
Modelo                    & Tokenizador  & \emph{Prompts}     & Acc       & AUROC     & F1-macro \\
\hline
DermLIP v1/v2             & PubMedBERT   & Genéricos          & $0{,}086$ & $0{,}366$ & $0{,}023$ \\
DermLIP v1/v2             & PubMedBERT   & Derm1M A3          & ---       & $0{,}516$ & --- \\
DermLIP original          & GPT-2/CLIP   & Derm1M A3          & $\mathbf{0{,}427}$ & $\mathbf{0{,}854}$ & --- \\
\hline
\end{tabular}
\end{table}

El rendimiento depende drásticamente de cómo se construye el texto.
Pasar de \emph{prompts} genéricos en inglés con tokenizador
PubMedBERT al conjunto A3 de Derm1M con tokenizador GPT-2/CLIP eleva
la AUROC de $0{,}366$ a $0{,}854$ ($+48{,}8$\,pp), como ilustra la
figura~\ref{fig:zs-prompts}. Los
\emph{prompts} genéricos con PubMedBERT quedan en AUROC$<0{,}5$,
señal de que proyectan en un espacio no alineado con el encoder
visual de la variante v1/v2. Cambiar solo el conjunto A3, sin tocar
el tokenizador, sube la AUROC a $0{,}516$; cambiar solo el
tokenizador, manteniendo A3, la lleva a $0{,}854$. El principal cuello
de botella del régimen \emph{zero-shot} no reside, por tanto, en la
representación visual ---idéntica en todas las configuraciones---,
sino en la alineación entre el espacio textual y el espacio visual.

Esta diferencia de casi $50$\,pp de AUROC, obtenida sin modificar el
encoder visual ni las imágenes, tiene una implicación metodológica
directa: evaluar un modelo \emph{zero-shot} con una única configuración
de \emph{prompts} o tokenizadores puede falsear sus capacidades
reales, y una comparación entre modelos puede incluso invertirse según
el \emph{prompt} elegido.

La generalización de este comportamiento a escenarios clínicos
externos se evalúa sobre la tarea binaria melanoma/no-melanoma en
cuatro datasets independientes
(DDI, MSKCC, HIBA y Derm7pt clínico) con la configuración
PubMedBERT de DermLIP v1/v2 (tabla~\ref{tab:zs-binaria}).

\begin{table}[H]
\centering
\footnotesize
\setlength{\tabcolsep}{4pt}
\caption{Zero-shot binario melanoma/no-melanoma con DermLIP v1/v2
(\emph{prompts} genéricos).}
\label{tab:zs-binaria}
\begin{tabular}{lrrrrr}
\hline
Dataset           & $N_{\mathrm{test}}$ & Acc       & BAcc      & AUROC     & F1-macro \\
\hline
DDI               & $137$               & $0{,}766$ & $0{,}499$ & $0{,}436$ & $0{,}463$ \\
MSKCC             & $1\,664$            & $0{,}716$ & $0{,}500$ & $0{,}430$ & $0{,}417$ \\
HIBA              & $334$               & $0{,}859$ & $0{,}497$ & $0{,}502$ & $0{,}462$ \\
Derm7pt clín.     & $168$               & $0{,}589$ & $0{,}472$ & $0{,}337$ & $0{,}439$ \\
\hline
\end{tabular}
\end{table}

La sensibilidad observada en HAM10000 se reproduce en los cuatro
datasets externos: bajo la configuración PubMedBERT con \emph{prompts}
genéricos, DermLIP opera cerca del azar (BAcc $\approx 0{,}5$), lo que
confirma que el rendimiento \emph{zero-shot} depende críticamente de la
alineación texto-imagen y limita la robustez del paradigma zero-shot
puro (análisis extendido en el Anexo~\ref{cap:anexog}).

\begin{figure}[H]
\centering
\includegraphics[width=0.78\textwidth]{fig_zs_prompts.png}
\caption{Sensibilidad zero-shot al tokenizador y a la
formulación del \emph{prompt} sobre HAM10000. La AUROC oscila
entre $0{,}366$ (DermLIP v1/v2 + PubMedBERT + prompts
genéricos) y $0{,}854$ (DermLIP original +
GPT-2/CLIP + prompts A3 de Derm1M), un rango de
$48{,}8$\,pp para variaciones que no afectan al encoder
visual.}
\label{fig:zs-prompts}
\end{figure}

\section{Comparación con modelos de lenguaje multimodales}
\label{sec:llm-results}

Esta sección contrasta si los grandes modelos de lenguaje
multimodales, entrenados sobre datasets mucho mayores, cierran la
brecha frente a los especialistas. La
tabla~\ref{tab:llm-comparativa} evalúa seis modelos de lenguaje
multimodales sobre HAM10000, PAD-UFES-20 y DDI, junto a los modelos de
referencia de las secciones~\ref{sec:rep-panderm}
y~\ref{sec:ft-results}.

\begin{table}[H]
\centering
\footnotesize
\setlength{\tabcolsep}{4pt}
\renewcommand{\arraystretch}{0.95}
\caption{Comparación de paradigmas sobre HAM10000, PAD-UFES-20 y DDI
($N_{\mathrm{test}} = 1\,232 / 461 / 137$). FT: ajuste fino, LP:
sondeo lineal, ZS: zero-shot.}
\label{tab:llm-comparativa}
\begin{tabular}{llrrrrrr}
\hline
Modelo                 & Modo  & HAM Acc       & HAM BAcc      & PAD Acc       & PAD BAcc      & DDI Acc       & DDI BAcc \\
\hline
PanDerm Base (TTA HAM) & FT    & $\mathbf{0{,}920}$ & $\mathbf{0{,}852}$ & $0{,}755$     & $\mathbf{0{,}704}$ & $0{,}774$     & $0{,}693$ \\
PanDerm Large          & LP    & $0{,}888$     & $0{,}575$     & $\mathbf{0{,}772}$ & $0{,}642$     & $\mathbf{0{,}847}$ & $\mathbf{0{,}764}$ \\
SigLIP-Large           & LP    & $0{,}900$     & ---           & $0{,}718$     & ---           & $0{,}789$     & --- \\
MedGemma 27B + LoRA    & FT    & $0{,}802$     & $0{,}270$     & $0{,}430$     & $0{,}347$     & $0{,}657$     & $0{,}618$ \\
MedGemma 27B           & ZS    & $0{,}665$     & $0{,}243$     & $0{,}447$     & $0{,}342$     & $0{,}474$     & $0{,}566$ \\
GPT-4o                 & ZS    & $0{,}485$     & $0{,}465$     & $0{,}553$     & $0{,}523$     & $0{,}723$     & $0{,}648$ \\
GPT-4o-mini            & ZS    & $0{,}256$     & $0{,}319$     & $0{,}390$     & $0{,}435$     & $0{,}737$     & $0{,}556$ \\
Gemini 2.5 Pro         & ZS    & $0{,}403$     & ---           & $0{,}477$     & ---           & $0{,}441$     & --- \\
Gemini 2.5 Flash       & ZS    & $0{,}380$     & ---           & $0{,}486$     & ---           & $0{,}555$     & --- \\
DermLIP                & ZS    & $0{,}239$     & $0{,}581$     & $0{,}601$     & ---           & $0{,}329$     & --- \\
BLIP-2 (Flan-T5-XL)    & ZS    & $0{,}015$     & $0{,}145$     & $0{,}017$     & $0{,}167$     & $0{,}796$     & $0{,}542$ \\
InstructBLIP (Flan-T5-XL)& ZS  & $0{,}024$     & $0{,}123$     & $0{,}171$     & $0{,}224$     & $0{,}788$     & $0{,}500$ \\
\hline
\end{tabular}
\end{table}

Los resultados separan claramente especialistas y generalistas. El
mejor sistema dermatológico (PanDerm Base ajustado con TTA) alcanza
accuracy $0{,}920$ sobre HAM10000, frente a $0{,}485$ del mejor
generalista \emph{zero-shot} (GPT-4o) ---una diferencia de $43{,}5$\,pp
que se mantiene en PAD-UFES-20 y DDI
(figura~\ref{fig:llm-comparison})---; en el extremo inferior, BLIP-2 e
InstructBLIP caen a accuracy $\sim 0{,}02$ sobre HAM10000. La lectura
no es que un modelo generalista sea deficiente, sino que, cuando el
problema es extremadamente visual y especializado, el conocimiento
dermatológico adquirido durante el preentrenamiento domina sobre la
mayor capacidad lingüística de los modelos multimodales
generalistas.\footnote{Las cifras de los modelos accesibles por API
comercial corresponden al \emph{prompt} y a la versión del proveedor
disponibles en la fecha de evaluación, no reproducible ante
actualizaciones silenciosas (sección~\ref{sec:lim-pesos}); documentan
un punto de comparación reproducible en su formulación, no el techo
intrínseco de cada sistema.}

\subsection{Ajuste supervisado de un LLM multimodal: MedGemma con LoRA}

MedGemma es el único LLM multimodal que se acerca, y lo hace gracias a
la adaptación al dominio: el ajuste con LoRA sobre MedGemma~27B mejora
$+13{,}7$\,pp de accuracy sobre HAM10000 respecto a su variante
\emph{zero-shot} ($0{,}665 \to 0{,}802$) ---un salto que vuelve a
señalar que es la adaptación, y no la escala bruta, lo que cierra la
brecha---. Aun así, incluso adaptado queda $11{,}8$\,pp por debajo del
mejor especialista y la mejora no transfiere de forma sistemática
(sin ganancia sobre PAD-UFES-20). El desglose por dataset, los
patrones de error de GPT-4o y los costes de inferencia figuran en el
Anexo~\ref{cap:anexog}.

\begin{figure}[H]
\centering
\includegraphics[width=0.92\textwidth]{fig_llm_comparison.png}
\caption{Comparativa de paradigmas sobre HAM10000. Gradiente
de $43{,}5$\,pp accuracy entre el mejor especialista
ajustado (PanDerm Base FT con TTA, $0{,}920$) y el mejor LLM
multimodal generalista \emph{zero-shot} (GPT-4o,
$0{,}485$). BLIP-2 e InstructBLIP descienden hasta accuracy
$\sim 0{,}02$.}
\label{fig:llm-comparison}
\end{figure}

\section{Equidad por fototipo cutáneo}
\label{sec:equidad-results}

Esta sección analiza si los modelos rinden de forma equitativa entre
fototipos cutáneos y cómo distinguir una equidad real de una métrica
engañosa: un \emph{gap} pequeño puede esconder un modelo
uniformemente peor, de modo que la métrica aislada no decide, como
establecen la comparación de cinco encoders y el test de DeLong. El rendimiento global de los cinco encoders sobre Fitzpatrick17k
completo ($N = 16\,577$), en las dos formulaciones consideradas
---$114$ patologías y $3$ clases de malignidad---, figura en el
Anexo~\ref{cap:anexob} (tabla~\ref{tab:equidad-global}); el análisis
de equidad se centra en la formulación de $3$ clases. La
tabla~\ref{tab:equidad-fototipo} reporta la BAcc por fototipo
Fitzpatrick sobre esa formulación.

\begin{table}[H]
\centering
\footnotesize
\setlength{\tabcolsep}{4pt}
\renewcommand{\arraystretch}{0.95}
\caption{BAcc por subgrupo Fitzpatrick (3 clases de malignidad) y
\emph{gap} entre fototipos extremos. Tamaños muestrales por
subgrupo: $299 / 470 / 325 / 261 / 169 / 73$ para FST~I a~VI.}
\label{tab:equidad-fototipo}
\begin{tabular}{lrrrrrrr}
\hline
Modelo                    & FST I         & FST II        & FST III       & FST IV        & FST V         & FST VI        & Gap I-VI \\
\hline
PanDerm Large             & $0{,}723$     & $0{,}708$     & $0{,}743$     & $0{,}669$     & $0{,}568$     & $0{,}506$     & $+0{,}217$ \\
PanDerm Base              & $0{,}653$     & $0{,}693$     & $0{,}669$     & $0{,}630$     & $0{,}585$     & $0{,}384$     & $+0{,}269$ \\
DermLIP v2 (visual)       & $0{,}716$     & $0{,}730$     & $0{,}765$     & $0{,}710$     & $0{,}693$     & $0{,}434$     & $+0{,}281$ \\
DINOv2 ViT-L/14           & $0{,}673$     & $0{,}671$     & $0{,}679$     & $0{,}654$     & $0{,}553$     & $0{,}368$     & $+0{,}306$ \\
BiomedCLIP                & $0{,}569$     & $0{,}602$     & $0{,}606$     & $0{,}566$     & $0{,}590$     & $\mathbf{0{,}445}$ & $\mathbf{+0{,}124}$ \\
\hline
\end{tabular}
\end{table}

Los cinco modelos tienen \emph{gap}~I-VI positivo: todos rinden peor
sobre los fototipos elevados (figura~\ref{fig:fairness-5models}). Pero
el ranking de equidad no coincide con el de rendimiento. DermLIP v2
logra la mejor BAcc global ($0{,}721$ sobre $3$ clases) y la mejor
AUROC ($0{,}909$), pero con un \emph{gap}~I-VI ($+0{,}281$) mayor que
el de PanDerm Large ($+0{,}217$); y BiomedCLIP exhibe el \emph{gap}
menor ($+0{,}124$) pero la peor BAcc global ($0{,}590$).

Este contraste encierra la principal aportación metodológica de la
sección: \textbf{el \emph{gap} aislado no constituye una métrica
suficiente de equidad}. Un modelo uniformemente peor entre subgrupos
puede presentar una brecha reducida simplemente porque falla en todos
ellos ---ser equitativamente malo no es ser equitativo en sentido
sustantivo---. La interpretación debe además acotarse por el reducido
tamaño muestral del fototipo~VI ($N = 73$, frente a $299$ en
FST~I; sección~\ref{sec:lim-muestrales}), que ensancha los intervalos
del \emph{gap} sobre los fototipos oscuros. El contraste con cierre
estadístico se reserva al test de DeLong
(sección~\ref{sec:equidad-delong}).

Un experimento cruzado \emph{Light}~$\leftrightarrow$~\emph{Dark}
(\emph{Light} agrupa FST~I-III; \emph{Dark}, FST~IV-VI; matriz completa
en el Anexo~\ref{cap:anexob}, tabla~\ref{tab:cross-eval}) añade una
segunda lectura, más fuerte que el simple \emph{gap}.
Entrenar solo con fototipos claros y evaluar sobre oscuros (L$\to$D
frente a L$\to$L) penaliza el rendimiento. La caída es mayor en
DINOv2 ($\Delta = -0{,}132$\,pp BAcc) y menor en BiomedCLIP
($\Delta = -0{,}063$\,pp). DermLIP v2 muestra la menor caída
Dark$\to$Light frente a Dark$\to$Dark ($\Delta = +0{,}016$\,pp), lo
que apunta a representaciones más invariantes al fototipo cutáneo
(figura~\ref{fig:cross-eval}).

\begin{figure}[H]
\centering
\includegraphics[width=0.95\textwidth]{figura_fairness_5models.png}
\caption{Equidad por fototipo Fitzpatrick (formulación de tres
clases de malignidad) sobre los cinco encoders evaluados. Los
cinco presentan \emph{gap}~I-VI positivo (entre $+0{,}124$ y
$+0{,}306$\,pp BAcc). Paradoja BiomedCLIP: \emph{gap} mínimo
asociado a peor BAcc global, lo que demuestra el carácter no
decisional de la métrica aislada.}
\label{fig:fairness-5models}
\end{figure}

\begin{figure}[H]
\centering
\includegraphics[width=0.85\textwidth]{fig_cross_eval.png}
\caption{Evaluación cruzada
\emph{Light}~$\leftrightarrow$~\emph{Dark} sobre Fitzpatrick17k
(BAcc en formulación de tres clases). La caída
L$\to$D vs L$\to$L cuantifica la pérdida de rendimiento al
entrenar sólo con fototipos claros y evaluar sobre fototipos
oscuros, máxima en DINOv2 ($\Delta = -0{,}132$\,pp).}
\label{fig:cross-eval}
\end{figure}

La proximidad entre algunos resultados ---especialmente entre PanDerm
Large y DermLIP v2--- plantea si las diferencias reflejan una ventaja
real o son compatibles con la variabilidad muestral, lo que motiva un
contraste formal mediante el test de DeLong.

\subsection{Validación estadística del ranking de encoders: test de
DeLong}
\label{sec:equidad-delong}

Las diferencias de AUROC entre encoders se contrastan mediante el test
de DeLong~\cite{delong1988}, que controla la correlación de evaluar los
cinco modelos sobre las mismas imágenes, sobre el endpoint binario de
malignidad (\emph{maligno} frente al resto) en el test de
Fitzpatrick17k ($N = 1\,658$, $226$ malignos), con regresión logística
binaria por encoder sobre sus \emph{features} congeladas
(sección~\ref{sec:lp}). La tabla~\ref{tab:delong-auroc} reporta la
AUROC binaria con IC95\% por la varianza de DeLong; la matriz completa
de $p$-valores pareados (dos colas) figura en el
Anexo~\ref{cap:anexob} (tabla~\ref{tab:delong-pares}).

\begin{table}[H]
\centering
\footnotesize
\setlength{\tabcolsep}{4pt}
\renewcommand{\arraystretch}{0.95}
\caption{AUROC binaria de detección de malignidad (\emph{maligno}
frente al resto) sobre el test de Fitzpatrick17k
($N = 1\,658$, $226$ malignos), con intervalo de confianza al
$95\%$ por la varianza de DeLong. Mejor valor en negrita.}
\label{tab:delong-auroc}
\begin{tabular}{lrl}
\hline
Modelo            & AUROC               & IC95\% (DeLong) \\
\hline
PanDerm Large     & $\mathbf{0{,}9488}$ & $[0{,}9350,\, 0{,}9626]$ \\
DermLIP v2        & $0{,}9452$          & $[0{,}9307,\, 0{,}9598]$ \\
PanDerm Base      & $0{,}9296$          & $[0{,}9130,\, 0{,}9463]$ \\
DINOv2 ViT-L/14   & $0{,}9150$          & $[0{,}8932,\, 0{,}9367]$ \\
BiomedCLIP        & $0{,}8915$          & $[0{,}8672,\, 0{,}9158]$ \\
\hline
\end{tabular}
\end{table}

El test de DeLong confirma que PanDerm Large y DermLIP v2 presentan un
rendimiento estadísticamente indistinguible
($\Delta\mathrm{AUROC} = 0{,}0036$, $p = 0{,}55$): ningún dato sostiene
una jerarquía entre ambos sobre este \emph{endpoint}. Los dos superan,
en cambio, de forma significativa a los encoders generalistas (DINOv2
y BiomedCLIP) y a PanDerm Base ($p \leq 7{,}3\times 10^{-3}$ en los
seis contrastes), lo que respalda cuantitativamente la ventaja del
preentrenamiento específico de dominio. Tras corrección de Bonferroni
(umbral $0{,}005$), la mayoría de estas diferencias se mantienen
significativas, mientras que algunas comparaciones entre generalistas
(DINOv2--BiomedCLIP, DINOv2--PanDerm~Base) dejan de ser concluyentes,
delimitando qué distinciones del \emph{ranking} son señal y cuáles
ruido muestral.

La lectura que el contraste habilita es más sólida que coronar a un
único modelo: los resultados no identifican un encoder dermatológico
claramente superior, sino una \emph{familia} de encoders
dermatológicos que supera de forma consistente a los modelos
generalistas ---conclusión que se retoma en la discusión
(sección~\ref{sec:disc-especializacion}).

\section{Validación en castellano sobre DermapixelAI 1.0}
\label{sec:dermapixel-results}

Los bloques anteriores evalúan los modelos sobre datasets externos
anglosajones. El objetivo de generalización de la
sección~\ref{sec:pregunta} exige comprobar si los hallazgos se
sostienen sobre material clínico en castellano, independiente y libre
de fuga frente al preentrenamiento. DermapixelAI~1.0 reúne ambas
condiciones. El bloque que se extiende hasta la
sección~\ref{sec:spanderm-clip} responde a si los modelos
fundacionales transfieren a este material y comparte un hallazgo
transversal: no existe un método óptimo único, sino que la mejor
configuración depende del nivel ontológico evaluado. La
tabla~\ref{tab:dermapixel-taxonomia} fija la nomenclatura de los
cuatro sistemas contrastados.

\begin{table}[H]
\centering
\footnotesize
\setlength{\tabcolsep}{4pt}
\renewcommand{\arraystretch}{0.95}
\caption{Nomenclatura de los cuatro sistemas evaluados sobre
DermapixelAI~1.0 a lo largo del bloque
secciones~\ref{sec:dermapixel-results}--\ref{sec:spanderm-clip}.}
\label{tab:dermapixel-taxonomia}
\begin{tabular}{lll}
\hline
Sistema & Qué es & Sección \\
\hline
PanDerm LP                          & Codificador congelado $+$ regresión logística           & §\ref{sec:dermapixel-results} \\
DermLIP \emph{zero-shot}            & Modelo visión-lenguaje externo (sin cabeza entrenada)   & §\ref{sec:dermapixel-zs} \\
Dermapixel R0 (cabeza superv.\ L2) & Rama supervisada con adaptación LoRA $r=16$             & §\ref{sec:dermapixel-spanderm-v0} \\
Dermapixel R0 (rama contrastiva)   & Rama \emph{CLIP-style} castellana \emph{zero-shot open} & §\ref{sec:spanderm-clip} \\
\hline
\end{tabular}
\end{table}

Frente al solapamiento de Dermnet
(sección~\ref{sec:auditoria-resultados}), DermapixelAI~1.0
(Anexo~\ref{cap:anexoe}) es completamente independiente del
preentrenamiento de la familia PanDerm/DermLIP: la auditoría por hash
MD5 sobre $507\,657$ imágenes catalogadas (Derm1M $415\,451$ ítems más
los once datasets externos del capítulo) reporta intersección nula
($0$ coincidencias sobre las $1\,062$ imágenes evaluables). Esta
evaluación no está, por tanto, sujeta a fuga de información.

\subsection{Transferencia congelada: sondeo lineal y validación cruzada}

\paragraph{Sondeo lineal sobre el \emph{split} fijo.}
Se aplica el protocolo de sondeo lineal (sección~\ref{sec:lp}) sobre
las $1\,062$ imágenes con \texttt{label\_source=ontology} (train
$N = 874$, validación $N = 152$, test $N = 36$), con PanDerm Base
($768$-D) y Large ($1\,024$-D) como extractores congelados y cabeza
\emph{LogisticRegression}. La evaluación L2 usa las $38$ subcategorías
efectivas tras consolidar las dos variantes ortográficas de
\emph{Trastornos de la queratinización} (Anexo~\ref{cap:anexoc},
sección~\ref{sec:anexo-l2-l3}).

\begin{table}[H]
\centering
\footnotesize
\setlength{\tabcolsep}{4pt}
\renewcommand{\arraystretch}{0.95}
\caption{Sondeo lineal con PanDerm Base y Large sobre
DermapixelAI~1.0 en los tres niveles de la ontología jerárquica
($N_{\mathrm{train}} = 874$, $N_{\mathrm{test}} = 36$ en L1/L2 y
$N_{\mathrm{test}} = 28$ en L3 tras filtrar clases ausentes en
train). Entre corchetes, IC95\% por \emph{bootstrap} estratificado
($1\,000$ remuestreos). Mejor valor por columna y nivel en
negrita.}
\label{tab:dermapixel-lp}
\adjustbox{max width=\textwidth}{%
\begin{tabular}{llrrrrr}
\hline
Nivel & Modelo & Acc & BAcc & AUROC & W-F1 & Kappa \\
\hline
L1 (4 clases)   & Base  & $0{,}639$ [$0{,}472$--$0{,}778$]
                        & $0{,}570$ [$0{,}371$--$0{,}751$]
                        & $0{,}697$ [$0{,}527$--$0{,}841$]
                        & $0{,}624$ [$0{,}457$--$0{,}760$]
                        & $0{,}296$ \\
L1 (4 clases)   & Large & $\mathbf{0{,}750}$ [$0{,}639$--$0{,}861$]
                        & $\mathbf{0{,}735}$ [$0{,}637$--$0{,}835$]
                        & $\mathbf{0{,}796}$ [$0{,}656$--$0{,}930$]
                        & $\mathbf{0{,}722}$ [$0{,}596$--$0{,}845$]
                        & $\mathbf{0{,}520}$ \\
L2 (38 clases)  & Base  & $0{,}222$ [$0{,}139$--$0{,}333$]
                        & $0{,}221$ [$0{,}147$--$0{,}294$]
                        & $0{,}707$ [$0{,}666$--$0{,}747$]
                        & $0{,}212$ [$0{,}092$--$0{,}302$]
                        & $0{,}170$ \\
L2 (38 clases)  & Large & $\mathbf{0{,}250}$ [$0{,}167$--$0{,}333$]
                        & $\mathbf{0{,}265}$ [$0{,}206$--$0{,}338$]
                        & $\mathbf{0{,}793}$ [$0{,}749$--$0{,}837$]
                        & $\mathbf{0{,}223}$ [$0{,}139$--$0{,}291$]
                        & $\mathbf{0{,}199}$ \\
L3 (224 clases) & Base  & $0{,}143$ [$0{,}107$--$0{,}179$]
                        & $0{,}132$ [$0{,}105$--$0{,}158$]
                        & $0{,}735$ [$0{,}693$--$0{,}779$]
                        & $0{,}107$ [$0{,}064$--$0{,}141$]
                        & $0{,}112$ \\
L3 (224 clases) & Large & $\mathbf{0{,}179}$ [$0{,}143$--$0{,}214$]
                        & $\mathbf{0{,}184}$ [$0{,}158$--$0{,}211$]
                        & $\mathbf{0{,}813}$ [$0{,}769$--$0{,}856$]
                        & $\mathbf{0{,}143}$ [$0{,}100$--$0{,}184$]
                        & $\mathbf{0{,}154}$ \\
\hline
\end{tabular}}
\end{table}

PanDerm Large supera a Base en las cinco métricas y los tres niveles
(entre $+7{,}8$ y $+9{,}9$\,pp de AUROC y entre $+4{,}4$ y
$+16{,}5$\,pp de BAcc), y el rendimiento decrece monótonamente con la
granularidad: la BAcc de Large cae de $0{,}735$ (L1) a $0{,}265$ (L2)
y $0{,}184$ (L3, $224$ diagnósticos), si bien aún multiplica el azar
por $2{,}9\times / 10{,}2\times / 40{,}9\times$. La cifra L1
BAcc $= 0{,}735$ (IC95\% [$0{,}637$--$0{,}835$]) se sitúa dentro del
rango de los diez datasets externos (tabla~\ref{tab:lp-panderm}), lo
que ubica a DermapixelAI~1.0 como \emph{benchmark} exigente en
castellano sin marcas dermatoscópicas estandarizadas.


\paragraph{Validación cruzada agrupada por caso.}
Dado que el test fijo $N = 36$ es un único \emph{split}, se estima el
rendimiento esperado mediante validación cruzada de cinco \emph{folds}
agrupada por \texttt{case\_id} (\emph{StratifiedGroupKFold} por L1
sobre las $1\,062$ imágenes), sin que ninguna imagen del mismo caso
caiga a la vez en \emph{train} y \emph{test}. La
tabla~\ref{tab:dermapixel-cv} reporta media y desviación sobre los
cinco \emph{folds}.

\begin{table}[H]
\centering
\footnotesize
\setlength{\tabcolsep}{4pt}
\renewcommand{\arraystretch}{0.95}
\caption{Validación cruzada agrupada por caso (cinco \emph{folds},
estratificado por L1 sobre \texttt{case\_id}) con PanDerm Large
sobre DermapixelAI~1.0. Cifras: media $\pm$ desviación estándar
sobre los cinco \emph{folds}. La cardinalidad L3 reportada aquí
(250 clases) corresponde a las clases efectivas en el dataset
completo; las restantes tablas del split fijo
(tabla~\ref{tab:dermapixel-lp}, tabla~\ref{tab:dermapixel-ensemble}
y las del Anexo~\ref{cap:anexoj}) usan 224 clases L3, las
únicas representadas en el split \texttt{train}.}
\label{tab:dermapixel-cv}
\begin{tabular}{lrrrrrr}
\hline
Nivel           & Acc@1                 & Acc@3                 & BAcc                  & AUROC                 & W-F1                  & Kappa \\
\hline
L1 (4 cls)      & $0{,}574 \pm 0{,}029$ & $0{,}994 \pm 0{,}004$ & $0{,}440 \pm 0{,}058$ & $0{,}753 \pm 0{,}054$ & $0{,}569 \pm 0{,}025$ & $0{,}314 \pm 0{,}035$ \\
L2 (38 cls)     & $0{,}227 \pm 0{,}015$ & $0{,}418 \pm 0{,}037$ & $0{,}168 \pm 0{,}020$ & $0{,}749 \pm 0{,}019$ & $0{,}221 \pm 0{,}022$ & $0{,}186 \pm 0{,}015$ \\
L3 (250 cls)    & $0{,}124 \pm 0{,}007$ & $0{,}232 \pm 0{,}028$ & $0{,}108 \pm 0{,}014$ & $0{,}800 \pm 0{,}018$ & $0{,}124 \pm 0{,}005$ & $0{,}110 \pm 0{,}005$ \\
\hline
\end{tabular}
\end{table}

La validación cruzada agrupada por caso es la estimación defendible
del rendimiento esperado. El test fijo sobrestima la media de los
cinco \emph{folds} en $+29{,}5$\,pp BAcc en L1, $+9{,}7$\,pp en L2 y
$+7{,}6$\,pp en L3, por la distribución favorable del split publicado
(sin \emph{Genodermatosis} en test y baja densidad de cola larga L3):
el split fijo es, por tanto, optimista. La AUROC, en cambio, se
mantiene estable frente a la partición (caídas $<5$\,pp:
$0{,}796 \to 0{,}753$ en L1, $0{,}793 \to 0{,}749$ en L2,
$0{,}813 \to 0{,}800$ en L3). Las cifras del test fijo
(tabla~\ref{tab:dermapixel-lp}) se conservan como punto operativo y
base de comparación con el experimento \emph{zero-shot}
(sección~\ref{sec:dermapixel-zs}), que reutiliza el mismo test.


\subsection{Validación del procedimiento de etiquetado}

Dado que parte de las etiquetas etiológicas procede de la ontología y
no de revisión individual, se compara el sondeo lineal sobre el
subconjunto verificado (\texttt{rosa\_verified}) frente al dataset
completo. Las diferencias son pequeñas, inconsistentes en dirección y
compatibles con la variabilidad muestral: la imputación etiológica no
degrada apreciablemente el rendimiento (resultados completos en el
Anexo~\ref{cap:anexoj}, tabla~\ref{tab:dermapixel-rosa}). Diversas
ablaciones adicionales sobre los mismos \emph{embeddings} (cabezas
$k$-NN/MLP, función de pérdida, \emph{Test-Time Augmentation}) tampoco
modifican las conclusiones y se documentan en el mismo anexo.


\subsection{Estrategias de adaptación más allá del sondeo lineal}
\label{sec:dermapixel-ensemble}

Más allá de la representación congelada, dos estrategias buscan mejorar
el rendimiento sobre DermapixelAI~1.0: combinar codificadores
complementarios y ajustar finamente el encoder.

\paragraph{Combinación de codificadores.}
Para evaluar si las representaciones de distintos codificadores son
complementarias, se entrena una cabeza logística por codificador
(PanDerm Base, Large y DermLIP v2; protocolo de sección~\ref{sec:lp})
sobre los \emph{embeddings} congelados y se contrastan frente a cinco
estrategias de combinación (promedio de las tres, promedio
Large$+$DermLIP, máximo por clase y apilado con meta-cabeza). La
tabla~\ref{tab:dermapixel-ensemble} reporta los tres niveles en el
mismo test fijo.

\begin{table}[H]
\centering
\footnotesize
\setlength{\tabcolsep}{4pt}
\renewcommand{\arraystretch}{0.95}
\caption{Combinación de codificadores sobre DermapixelAI~1.0. LP
individual sobre PanDerm Base, PanDerm Large y DermLIP v2; tres
estrategias de combinación: promedio Large$+$DermLIP, máximo por
clase entre los tres, y apilado con meta-cabeza. Mejor valor por
columna y nivel en negrita.}
\label{tab:dermapixel-ensemble}
\begin{tabular}{llrrrr}
\hline
Nivel & Estrategia                         & Acc@1              & Acc@3              & BAcc               & AUROC \\
\hline
L1 (4 cls)   & LP PanDerm Base             & $0{,}639$          & ---                & $0{,}570$          & $0{,}697$ \\
L1 (4 cls)   & LP PanDerm Large            & $\mathbf{0{,}750}$ & $\mathbf{1{,}000}$ & $\mathbf{0{,}735}$ & $0{,}796$ \\
L1 (4 cls)   & LP DermLIP v2               & $0{,}694$          & $\mathbf{1{,}000}$ & $0{,}680$          & $0{,}789$ \\
L1 (4 cls)   & Avg L$+$D                   & $\mathbf{0{,}750}$ & $\mathbf{1{,}000}$ & $0{,}702$          & $\mathbf{0{,}817}$ \\
L1 (4 cls)   & Avg 3                       & $0{,}722$          & $\mathbf{1{,}000}$ & $0{,}694$          & $0{,}811$ \\
L1 (4 cls)   & Máx 3                       & $0{,}694$          & $\mathbf{1{,}000}$ & $0{,}652$          & $0{,}805$ \\
L1 (4 cls)   & Stack 3                     & $0{,}694$          & $\mathbf{1{,}000}$ & $0{,}703$          & $0{,}790$ \\
\hline
L2 (38 cls)  & LP PanDerm Base             & $0{,}222$          & ---                & $0{,}221$          & $0{,}707$ \\
L2 (38 cls)  & LP PanDerm Large            & $0{,}250$          & $0{,}500$          & $0{,}265$          & $0{,}793$ \\
L2 (38 cls)  & LP DermLIP v2               & $\mathbf{0{,}333}$ & $0{,}528$          & $\mathbf{0{,}279}$ & $\mathbf{0{,}860}$ \\
L2 (38 cls)  & Avg L$+$D                   & $0{,}278$          & $\mathbf{0{,}556}$ & $\mathbf{0{,}279}$ & $0{,}832$ \\
L2 (38 cls)  & Avg 3                       & $0{,}278$          & $0{,}528$          & $0{,}268$          & $0{,}821$ \\
L2 (38 cls)  & Máx 3                       & $0{,}222$          & $0{,}528$          & $0{,}235$          & $0{,}797$ \\
L2 (38 cls)  & Stack 3                     & $0{,}250$          & $0{,}528$          & $0{,}258$          & $0{,}804$ \\
\hline
L3 (224 cls) & LP PanDerm Base             & $0{,}143$          & ---                & $0{,}132$          & $0{,}735$ \\
L3 (224 cls) & LP PanDerm Large            & $\mathbf{0{,}179}$ & $\mathbf{0{,}286}$ & $\mathbf{0{,}184}$ & $\mathbf{0{,}813}$ \\
L3 (224 cls) & LP DermLIP v2               & $0{,}107$          & $0{,}214$          & $0{,}079$          & $0{,}809$ \\
L3 (224 cls) & Avg L$+$D                   & $\mathbf{0{,}179}$ & $0{,}250$          & $\mathbf{0{,}184}$ & $0{,}802$ \\
L3 (224 cls) & Avg 3                       & $0{,}143$          & $0{,}250$          & $0{,}152$          & $0{,}795$ \\
L3 (224 cls) & Máx 3                       & $0{,}143$          & $0{,}250$          & $0{,}132$          & $0{,}780$ \\
L3 (224 cls) & Stack 3                     & $0{,}143$          & $0{,}250$          & $0{,}158$          & $0{,}791$ \\
\hline
\end{tabular}
\end{table}

Ningún codificador lidera simultáneamente L1, L2 y L3. En L1 lidera la
combinación avg Large$+$DermLIP en AUROC ($0{,}796 \to 0{,}817$); en
L2, DermLIP v2 individual en todas las métricas ($+6{,}7$\,pp AUROC
sobre PanDerm Large, $0{,}793 \to 0{,}860$), sin que ningún
\emph{ensemble} lo supere, en coherencia con el experimento
\emph{zero-shot} (sección~\ref{sec:dermapixel-zs}); en L3 lidera
PanDerm Large individual, donde la cola larga limita la señal
complementaria.


\paragraph{Ajuste fino parcial del codificador.}
\label{sec:dermapixel-ft}
\label{sec:dermapixel-ft-l1}
Este protocolo replica sobre DermapixelAI~1.0 el ajuste fino
supervisado (FT) de la sección~\ref{sec:ft-results}. El encoder es
PanDerm Large con las dos últimas capas descongeladas ($25{,}2$\,M
parámetros entrenables, $8{,}3\%$ del encoder) y cabeza
\emph{fully-connected}, con entropía cruzada reponderada por el
desbalance de \emph{Genodermatosis} (receta en
Anexo~\ref{cap:anexoa}, tabla~\ref{tab:anexoa-ft-recipe}).

Sobre el nivel L1, el ajuste fino parcial alcanza BAcc $0{,}622$
(IC95\% [$0{,}510$--$0{,}735$]) y AUROC $0{,}715$, lo que mejora en
$+18{,}2$\,pp la cifra de validación cruzada agrupada por caso del
sondeo lineal congelado (BAcc $0{,}440 \pm 0{,}058$,
tabla~\ref{tab:dermapixel-cv}), con intervalos no solapados: el
ajuste fino aporta valor real más allá del ruido de partición. No
obstante, no supera a la combinación lineal de codificadores
congelados \texttt{avg\_large\_dermlip}
(sección~\ref{sec:dermapixel-ensemble}; BAcc $0{,}702$, AUROC
$0{,}817$). La replicación sobre L2 y L3 mantiene el mismo
protocolo, variando únicamente la cabeza supervisada. La
tabla~\ref{tab:dermapixel-ft-multilevel} compara las métricas
sobre los tres niveles.

\begin{table}[H]
\centering
\footnotesize
\setlength{\tabcolsep}{4pt}
\renewcommand{\arraystretch}{0.95}
\caption{Ajuste fino parcial del codificador (PanDerm Large, dos
últimas capas descongeladas) sobre los tres niveles ontológicos de
DermapixelAI~1.0. Test fijo $N_{\mathrm{test}} = 36$ en L1/L2 y
$N_{\mathrm{test}} = 28$ en L3. La cabeza L2 de este protocolo
dimensiona $1\,024 \to 37$: son las $37$ clases L2 efectivas en el
\emph{train} del ajuste fino, de las $38$ del corpus tras fusionar
las dos variantes ortográficas de ``queratinización'' ($39$ en
bruto); las $38$ clases son la cardinalidad nominal del nivel L2.
Entre corchetes, IC95\% por \emph{bootstrap}
estratificado ($1\,000$ remuestreos). Mejor valor por columna en
negrita.}
\label{tab:dermapixel-ft-multilevel}
\begin{tabular}{lrrrr}
\hline
Nivel        & Acc@1                            & Acc@3                            & BAcc                             & AUROC \\
\hline
L1 (4 cls)   & $\mathbf{0{,}583}$ [$0{,}444$--$0{,}722$] & $\mathbf{0{,}972}$ [$0{,}917$--$1{,}000$] & $\mathbf{0{,}622}$ [$0{,}510$--$0{,}735$] & $0{,}715$ [$0{,}603$--$0{,}830$] \\
L2 (37 cls)  & $0{,}250$ [$0{,}167$--$0{,}333$] & $0{,}500$ [$0{,}389$--$0{,}611$] & $0{,}324$ [$0{,}265$--$0{,}397$] & $\mathbf{0{,}852}$ [$0{,}814$--$0{,}888$] \\
L3 (224 cls) & $0{,}179$ [$0{,}143$--$0{,}214$] & $0{,}286$ [$0{,}214$--$0{,}357$] & $0{,}184$ [$0{,}158$--$0{,}211$] & $0{,}804$ [$0{,}773$--$0{,}835$] \\
\hline
\end{tabular}
\end{table}

El ajuste fino denso parcial mejora al sondeo lineal en L1 y L2, pero
no es la forma de adaptación más eficiente, como revela la comparación
directa con la adaptación de bajo rango (subsección siguiente). En
conjunto, ni la combinación de codificadores ni el ajuste fino
identifican una opción óptima universal: \textbf{el rendimiento
depende simultáneamente del nivel ontológico, de la estrategia de
adaptación y del criterio de evaluación}, de modo que la búsqueda de un
único codificador \emph{ganador} resulta una simplificación excesiva
del problema (lectura de diseño en
sección~\ref{sec:disc-especializacion}).

\subsection{LoRA frente a ajuste fino denso}

Sobre el nivel L2, el FT parcial denso ($25{,}2$\,M parámetros
entrenables) alcanza BAcc $0{,}324$ y AUROC $0{,}852$, cifras que la
adaptación de bajo rango (LoRA~\cite{lora}) del protocolo Dermapixel
R0 (sección~\ref{sec:dermapixel-spanderm-v0}) supera ---BAcc
$0{,}363 \pm 0{,}007$ y AUROC $0{,}855 \pm 0{,}006$--- movilizando
$48$ veces menos parámetros
(tabla~\ref{tab:dermapixel-ft-vs-lora}).

\begin{table}[H]
\centering
\footnotesize
\setlength{\tabcolsep}{4pt}
\renewcommand{\arraystretch}{0.95}
\caption{Comparativa sobre el nivel L2 de DermapixelAI~1.0 entre el
sondeo lineal congelado (referencia), el ajuste fino parcial denso
con dos últimas capas descongeladas y el protocolo Dermapixel R0
(cabeza supervisada L2) con
adaptación LoRA. Mismo conjunto de test ($N_{\mathrm{test}} = 36$,
$37$--$38$ clases según consolidación de variantes ortográficas
de queratinización descrita en
sección~\ref{sec:dermapixel-spanderm-v0}). Mejor valor por
columna en negrita.}
\label{tab:dermapixel-ft-vs-lora}
\begin{tabular}{lrrr}
\hline
Protocolo                       & BAcc L2                         & AUROC L2                       & Params entren. \\
\hline
LP test fijo (referencia, tabla~\ref{tab:dermapixel-lp}) & $0{,}265$                       & $0{,}793$                      & $39\,K$ \\
FT parcial denso (este trabajo) & $0{,}324$                       & $0{,}852$                      & $25{,}2$\,M \\
\textbf{Dermapixel R0 (LoRA)}   & $\mathbf{0{,}363 \pm 0{,}007}$  & $\mathbf{0{,}855 \pm 0{,}006}$ & $\mathbf{524\,K}$ \\
\hline
\end{tabular}
\end{table}

Bajo el régimen de datos disponible ($N_{\mathrm{train}} = 874$), el
cuello de botella no es la capacidad del modelo sino la eficiencia de
la adaptación: LoRA supera al FT denso parcial movilizando el $2\%$ de
sus parámetros ($524\,K$ frente a $25{,}2$\,M, $48\times$ menos), lo
que justifica la decisión arquitectónica de Dermapixel R0
(sección~\ref{sec:dermapixel-spanderm-v0}).

\subsection{Limitaciones del ajuste fino en la cola larga}

Sobre el nivel L3, el ajuste fino denso parcial no supera al sondeo
lineal congelado (BAcc $0{,}184$ en ambos): la cola larga severa del
nivel ($77\%$ de las clases con $\leq 5$ ejemplos) anula el beneficio
del ajuste fino, lo que reserva L3 al ranking por prototipos coseno de
Dermapixel R0 (tabla~\ref{tab:dermapixel-spanderm-l3}) o a la
ampliación del dataset vía banco hospitalario del HUSLL.

Tres \emph{caveats} acotan la interpretación de las cifras de esta
sección: el conjunto de test es reducido ($N = 36$ en L1/L2,
$N = 28$ en L3), lo que ensancha los IC95\%; la categoría etiológica
\emph{Genodermatosis} no aparece en el test de L1 ($0 / 36$), por lo
que las métricas reflejan una clasificación efectiva de tres clases;
y la cola larga de L3 limita el rendimiento absoluto alcanzable bajo
cualquier protocolo de ajuste denso sobre el dataset actual.

\section{Validación \emph{zero-shot} multimodal sobre DermapixelAI 1.0}
\label{sec:dermapixel-zs}

Esta sección evalúa hasta qué punto un modelo fundacional
dermatológico transfiere al dataset clínico en castellano sin entrenar
cabeza supervisada. La brecha entre el régimen \emph{zero-shot} y el
supervisado depende del nivel ontológico: amplia en L1 y L3, pero
prácticamente nula en L2, donde DermLIP v2 \emph{zero-shot} alcanza
una AUROC ($0{,}794$) indistinguible de la del sondeo lineal
supervisado de PanDerm Large ($0{,}793$).

El experimento opera sobre el mismo test que la
sección~\ref{sec:dermapixel-results} ($N_{\mathrm{test}} = 36$ en
L1/L2, $28$ en L3), sin cabeza clasificadora y con no solapamiento
garantizado ($0 / 507\,657$ coincidencias MD5). Se evalúan tres
modelos contrastivos ---DermLIP v2 (Derm1M), SigLIP-SO400M (WebLI) y
BiomedCLIP (PubMed-15M)--- por \emph{argmax} sobre la similitud coseno
con un \emph{ensemble} de tres plantillas por nivel, contrastando
\emph{prompt} en castellano frente a inglés; L3 se evalúa sólo en
castellano por falta de traducción inglesa validada para sus $224$
clases (tabla~\ref{tab:dermapixel-zs}).

\begin{table}[H]
\centering
\footnotesize
\setlength{\tabcolsep}{4pt}
\renewcommand{\arraystretch}{0.95}
\caption{Clasificación \emph{zero-shot} multimodal sobre el
conjunto de test de DermapixelAI~1.0 ($N_{\mathrm{test}} = 36$ en
L1/L2, $N_{\mathrm{test}} = 28$ en L3). Entre corchetes, IC95\%
por \emph{bootstrap} estratificado ($1\,000$ remuestreos) para
las dos métricas primarias (Acc@1 y AUROC). Las métricas BAcc y
Acc@3 se reportan como valor puntual por concisión; los IC95\%
correspondientes están disponibles en el material reproducible.
L3 se evalúa únicamente en castellano. Mejor valor por columna,
nivel e idioma en negrita.}
\label{tab:dermapixel-zs}
\begin{tabular}{lllrrrr}
\hline
Nivel & Modelo & Lang & Acc@1 & Acc@3 & BAcc & AUROC \\
\hline
L1 (4 cls)   & BiomedCLIP    & es & $\mathbf{0{,}389}$ [$0{,}250$--$0{,}528$]
                                  & $\mathbf{0{,}944}$
                                  & $0{,}375$
                                  & $0{,}679$ [$0{,}562$--$0{,}781$] \\
L1 (4 cls)   & BiomedCLIP    & en & $0{,}389$ [$0{,}250$--$0{,}528$]
                                  & $\mathbf{0{,}972}$
                                  & $0{,}324$
                                  & $0{,}639$ [$0{,}524$--$0{,}738$] \\
L1 (4 cls)   & SigLIP-SO400M & es & $0{,}222$ [$0{,}111$--$0{,}361$]
                                  & $0{,}917$
                                  & $0{,}178$
                                  & $0{,}650$ [$0{,}571$--$0{,}735$] \\
L1 (4 cls)   & SigLIP-SO400M & en & $0{,}333$ [$0{,}194$--$0{,}472$]
                                  & $0{,}889$
                                  & $0{,}292$
                                  & $0{,}663$ [$0{,}595$--$0{,}719$] \\
L1 (4 cls)   & DermLIP v2    & es & $0{,}361$ [$0{,}222$--$0{,}500$]
                                  & $0{,}861$
                                  & $0{,}379$
                                  & $0{,}699$ [$0{,}600$--$0{,}794$] \\
L1 (4 cls)   & DermLIP v2    & en & $\mathbf{0{,}500}$ [$0{,}333$--$0{,}667$]
                                  & $0{,}917$
                                  & $\mathbf{0{,}441}$
                                  & $\mathbf{0{,}709}$ [$0{,}565$--$0{,}841$] \\
\hline
L2 (38 cls)  & BiomedCLIP    & es & $0{,}111$ [$0{,}111$--$0{,}111$]
                                  & $0{,}278$
                                  & $0{,}118$
                                  & $0{,}647$ [$0{,}604$--$0{,}688$] \\
L2 (38 cls)  & BiomedCLIP    & en & $0{,}083$ [$0{,}056$--$0{,}111$]
                                  & $0{,}250$
                                  & $0{,}088$
                                  & $0{,}655$ [$0{,}601$--$0{,}713$] \\
L2 (38 cls)  & SigLIP-SO400M & es & $0{,}139$ [$0{,}139$--$0{,}139$]
                                  & $0{,}250$
                                  & $0{,}235$
                                  & $0{,}636$ [$0{,}589$--$0{,}680$] \\
L2 (38 cls)  & SigLIP-SO400M & en & $0{,}194$ [$0{,}111$--$0{,}278$]
                                  & $0{,}278$
                                  & $0{,}221$
                                  & $0{,}697$ [$0{,}640$--$0{,}747$] \\
L2 (38 cls)  & DermLIP v2    & es & $\mathbf{0{,}222}$ [$0{,}167$--$0{,}278$]
                                  & $0{,}333$
                                  & $0{,}265$
                                  & $0{,}762$ [$0{,}721$--$0{,}805$] \\
L2 (38 cls)  & DermLIP v2    & en & $\mathbf{0{,}250}$ [$0{,}194$--$0{,}306$]
                                  & $\mathbf{0{,}417}$
                                  & $\mathbf{0{,}294}$
                                  & $\mathbf{0{,}794}$ [$0{,}747$--$0{,}837$] \\
\hline
L3 (224 cls) & BiomedCLIP    & es & $0{,}107$ [$0{,}036$--$0{,}179$]
                                  & $0{,}214$
                                  & $0{,}079$
                                  & $0{,}703$ [$0{,}673$--$0{,}732$] \\
L3 (224 cls) & SigLIP-SO400M & es & $0{,}071$ [$0{,}000$--$0{,}143$]
                                  & $0{,}214$
                                  & $0{,}053$
                                  & $0{,}696$ [$0{,}667$--$0{,}723$] \\
L3 (224 cls) & DermLIP v2    & es & $\mathbf{0{,}107}$ [$0{,}071$--$0{,}143$]
                                  & $\mathbf{0{,}214}$
                                  & $\mathbf{0{,}079}$
                                  & $\mathbf{0{,}740}$ [$0{,}719$--$0{,}763$] \\
\hline
\end{tabular}
\end{table}

La evaluación \emph{zero-shot} sobre DermapixelAI~1.0 arroja tres
hallazgos principales: (1)~DermLIP v2 es el modelo más competitivo en
los tres niveles ontológicos; (2)~el idioma del \emph{prompt} afecta
significativamente a los modelos preentrenados mayoritariamente en
inglés; y (3)~la necesidad de supervisión depende del nivel
ontológico ---crítica en L1 y L3, pero mínima en L2---.

\textit{Primero}, DermLIP v2 obtiene el mejor o equivalente
rendimiento en las cinco métricas sobre los tres niveles (salvo
Acc@3 L1, liderada por BiomedCLIP), coherente con su preentrenamiento
sobre Derm1M, el único dataset específicamente dermatológico de los
tres~\cite{derm1m}.

\textit{Segundo}, traducir el \emph{prompt} del inglés al castellano
reduce la Acc@1 en L1 entre $11$ y $14$ puntos porcentuales en los
modelos anglófonos (DermLIP v2 $-13{,}9$\,pp, de $0{,}500$ a $0{,}361$;
SigLIP-SO400M $-11{,}1$\,pp), mientras que BiomedCLIP mantiene su
cifra en ambos idiomas. La penalización es consistente con el
predominio del inglés en el preentrenamiento ---aunque el diseño no
permite atribuirla causalmente a ese único factor--- y motiva
directamente la rama textual castellana de Dermapixel R0: sin
penalización lingüística, una rama en castellano no estaría
justificada.

\textit{Tercero}, la comparación con el sondeo lineal supervisado
(tabla~\ref{tab:dermapixel-lp}) cuantifica la brecha entre un sistema
entrenado y uno \emph{zero-shot}: $+29{,}4$\,pp BAcc en L1 y
$+10{,}5$\,pp en L3 a favor del LP. El comportamiento se invierte, en
cambio, en L2, y aquí aparece el resultado más llamativo de la
sección: DermLIP v2 \emph{zero-shot} en inglés (AUROC $0{,}794$) iguala
al mejor sondeo lineal supervisado (PanDerm Large, $0{,}793$) dentro de
los intervalos de confianza. Esto sugiere que, a nivel de
subcategorías dermatológicas, una parte sustancial del conocimiento
discriminativo ya está codificada en el espacio multimodal
preentrenado y puede recuperarse mediante texto, sin entrenamiento
adicional; el grano grueso L1 y la cola larga L3 sí requieren, en
cambio, clasificador supervisado.

La tabla~\ref{tab:dermapixel-zs-gaps} resume las brechas
cuantificadas en esta sección.

\begin{table}[H]
\centering
\footnotesize
\setlength{\tabcolsep}{4pt}
\renewcommand{\arraystretch}{0.95}
\caption{Brechas \emph{zero-shot} cuantificadas sobre DermapixelAI~1.0:
penalización por idioma (castellano frente a inglés) y distancia al
sondeo lineal supervisado (LP).}
\label{tab:dermapixel-zs-gaps}
\begin{tabular}{llr}
\hline
Contraste                                & Métrica & \emph{Gap} \\
\hline
DermLIP v2, L1, castellano vs.\ inglés   & Acc@1 & $-13{,}9$\,pp \\
SigLIP-SO400M, L1, castellano vs.\ inglés & Acc@1 & $-11{,}1$\,pp \\
DermLIP v2, L2, castellano vs.\ inglés   & AUROC & $-3{,}2$\,pp \\
LP supervisado vs.\ \emph{zero-shot}, L1 & BAcc  & $+29{,}4$\,pp \\
LP supervisado vs.\ \emph{zero-shot}, L3 & BAcc  & $+10{,}5$\,pp \\
\hline
\end{tabular}
\end{table}

\subsection*{Caveats}

Las cautelas sobre el tamaño del conjunto de test ($N = 36$ en
L1/L2, $N = 28$ en L3) y sobre la ausencia de solapamiento frente
a Derm1M ya enunciadas para el sondeo lineal
(sección~\ref{sec:dermapixel-results}) aplican también a esta
evaluación zero-shot. Dos \emph{caveats} adicionales son
específicos del régimen zero-shot. Primero, el \emph{ensemble} de
tres plantillas textuales por nivel constituye una formulación
mínima de \emph{prompt engineering}; la literatura sobre
\emph{zero-shot} multimodal emplea típicamente entre diez y ochenta
plantillas generadas con asistencia de LLMs, lo que sitúa las
cifras reportadas como cota inferior del rendimiento alcanzable.
Segundo, la evaluación L3 únicamente en castellano impide
cuantificar la penalización por idioma sobre la cola larga
diagnóstica, dado que el vocabulario ontológico no dispone de
traducción inglesa validada para sus $224$ clases.

\section{Ajuste fino con LoRA: Dermapixel R0 (cabeza supervisada L2)}
\label{sec:dermapixel-spanderm-v0}

% NOTA: el encargo pedía \subsubsection (numeración 4.x.y.z); estas
% secciones carecen de \subsection intermedia, por lo que
% \subsubsection numeraba 4.x.0.y. Se emplea \subsection (4.x.y),
% coherente con el resto del Cap4 y sin el nivel 0 espurio.
\subsection{Diseño de Dermapixel R0}

Frente a los experimentos anteriores, que evalúan modelos existentes
como extractores congelados o \emph{zero-shot}, esta sección aborda la
contribución propia: \textbf{Dermapixel R0}, una adaptación del modelo
fundacional dermatológico PanDerm Large sobre DermapixelAI~1.0 mediante
adaptación de bajo rango (\emph{Low-Rank Adaptation}, LoRA)~\cite{lora}.
Se trata de una adaptación de un modelo fundacional preexistente, no de
un preentrenamiento a escala. La denominación R0 designa una versión
base sobre un corpus reducido. Con $524$\,K parámetros entrenables
($0{,}17\%$ del total), Dermapixel R0 obtiene la mejor balanced accuracy
(BAcc) L2 documentada sobre el dataset. Se reportan la cabeza
supervisada L2 ($38$ clases efectivas tras consolidar las variantes
ortográficas de queratinización) y el ranking L3 por prototipos coseno
sobre las $250$ clases efectivas. La elección de L2 responde a su
cardinalidad operativa ($3{,}5$ muestras/clase en
$N_{\mathrm{train}} = 874$); la definición ontológica y las
cardinalidades L1/L2/L3 se detallan en
Métodos~(sección~\ref{sec:ontologia}) y en el Anexo~\ref{cap:anexoc}.
La configuración completa (LoRA $r=16$, optimización, augmentaciones,
semillas) figura en el Anexo~\ref{cap:anexoa}
(tabla~\ref{tab:dermapixel-spanderm-config}).

\subsection{Resultados principales L2}

Sobre el conjunto de test L2, Dermapixel R0 alcanza una BAcc de
$0{,}363 \pm 0{,}007$ y una AUROC de $0{,}855 \pm 0{,}006$ (media
$\pm$ desviación sobre tres semillas), con comportamiento estable
entre ellas y sobreajuste leve. La tabla~\ref{tab:dermapixel-spanderm-comp}
desglosa la comparación con el sondeo lineal y DermLIP v2; el análisis
\emph{best-checkpoint} frente a las últimas épocas figura en el
Anexo~\ref{cap:anexoi}.

La tabla~\ref{tab:dermapixel-spanderm-comp} sitúa Dermapixel R0 frente
a las referencias internas sobre el mismo test y nivel: sondeo lineal
congelado bajo validación cruzada agrupada por caso
(tabla~\ref{tab:dermapixel-cv}), sondeo lineal sobre el \emph{split}
fijo (tabla~\ref{tab:dermapixel-lp}) y DermLIP v2 individual
(sección~\ref{sec:dermapixel-ensemble}, mejor opción L2 sin cabeza
entrenada).

\begin{table}[H]
\centering
\footnotesize
\setlength{\tabcolsep}{4pt}
\renewcommand{\arraystretch}{0.95}
\caption{Comparativa sobre el nivel L2 de DermapixelAI~1.0
($38$ clases con queratinización consolidada,
$N_{\mathrm{test}} = 36$) entre Dermapixel R0 y las referencias
internas del trabajo. La fila LP CV se considera la cifra honesta
del sondeo lineal congelado por su independencia del \emph{split}
fijo y por su anclaje en validación cruzada agrupada por caso
(sección~\ref{sec:dermapixel-results}). Mejor valor por
columna en negrita.}
\label{tab:dermapixel-spanderm-comp}
\adjustbox{max width=\textwidth}{%
\begin{tabular}{lllrr}
\hline
Método                                              & Encoder                       & Params entren.    & BAcc L2                         & AUROC L2 \\
\hline
LP test fijo (tabla~\ref{tab:dermapixel-lp})        & PanDerm Large $1\,024$ (cong.) & $39\,K$           & $0{,}265$                       & $0{,}793$ \\
LP $5$-fold CV (tabla~\ref{tab:dermapixel-cv})      & PanDerm Large $1\,024$ (cong.) & $39\,K$           & $0{,}168 \pm 0{,}020$           & $0{,}749 \pm 0{,}019$ \\
DermLIP v2 individual (sección~\ref{sec:dermapixel-ensemble}) & DermLIP $512$ (cong.)        & $19{,}5\,K$       & $0{,}279$                       & $\mathbf{0{,}860}$ \\
\textbf{Dermapixel R0} (este trabajo)               & PanDerm Large $1\,024$ $+$ LoRA & $524\,K$          & $\mathbf{0{,}363 \pm 0{,}007}$  & $0{,}855 \pm 0{,}006$ \\
\hline
\end{tabular}}
\end{table}

Dermapixel R0 iguala o supera a las tres referencias internas sobre
BAcc L2 ($0{,}363 \pm 0{,}007$ frente a $0{,}168 \pm 0{,}020$ del LP CV,
$0{,}265$ del LP fijo y $0{,}279$ de DermLIP v2 individual): es la mejor
BAcc L2 documentada sobre el dataset. Dado que la comparación se realiza
sobre el test fijo ($N_{\mathrm{test}} = 36$), la referencia defendible
del sondeo lineal congelado es la de validación cruzada agrupada por
caso (tabla~\ref{tab:dermapixel-cv}); la ventaja de R0 se mantiene
también frente a esa estimación más conservadora ($+19{,}5$\,pp BAcc).
Su AUROC L2 ($0{,}855 \pm 0{,}006$) empata estadísticamente con la de
DermLIP v2 nativo ($0{,}860$), de modo que la adaptación ligera de
PanDerm Large alcanza el mismo techo discriminativo que el
preentrenamiento contrastivo sobre Derm1M, con ventaja en las clases
minoritarias.

\subsection{Adaptación eficiente frente a ajuste denso}

La adaptación de bajo rango de Dermapixel R0 entrena $524$\,K
parámetros frente a los $25{,}2$\,M del ajuste fino parcial denso de la
sección~\ref{sec:dermapixel-ft-l1} (aproximadamente $48\times$ menos), y
obtiene mejor BAcc L2. La adaptación eficiente de parámetros (PEFT)
supera al ajuste fino parcial denso en este régimen de datos.

\subsection{Limitaciones en L3}

El nivel L3 se aborda mediante ranking por prototipos coseno, inviable
como cabeza supervisada directa bajo $3{,}5$ muestras/clase. El ranking
alcanza Acc@1 $= 0{,}167$ y Acc@5 $= 0{,}296$ ($\times 42$ y $\times 15$
sobre el azar), pero permanece lejos del régimen clínico: $178$ de $250$
clases con $\leq 5$ imágenes en \emph{train} producen prototipos
ruidosos. Las cifras completas figuran en el Anexo~\ref{cap:anexoj}
(sección~\ref{sec:anexoj-l3-prototipos},
tabla~\ref{tab:dermapixel-spanderm-l3}). La progresión hacia una cabeza
L3 directa queda condicionada a la incorporación del banco hospitalario
del HUSLL (umbral $\geq 30$ imágenes/clase L3,
sección~\ref{sec:dermapixel-futuro}).

\section{Dermapixel R0: rama contrastiva multimodal castellano
zero-shot sobre DermapixelAI~1.0}
\label{sec:spanderm-clip}

A diferencia de la cabeza supervisada L2, esta rama mapea texto
clínico libre en castellano sobre el espacio de imágenes sin
vocabulario cerrado, condición relevante para la búsqueda documental
sobre el archivo clínico. La progresión de cinco iteraciones
(v1$\to$v5) sobre una adaptación de PanDerm Large ---no un
preentrenamiento a escala--- indica que la palanca decisiva es la
capacidad del codificador visual, no la escala de datos externos, e
incluye un resultado negativo de escala. La condicionalidad del
\emph{zero-shot} L3 al volumen de pares y a la cobertura de la cola
larga se mantiene como cautela transversal.

\subsection{Diseño experimental}

La rama contrastiva sigue el patrón dual-encoder estilo
CLIP~\cite{clip2021}: un encoder visual basado en PanDerm Large y un
encoder textual multilingüe adaptado al lenguaje clínico castellano,
en sustitución del encoder anglosajón de DermLIP (la justificación
frente a un \emph{wrapper} inglés con traducción ES--EN, en el
Anexo~\ref{cap:anexoi}). Se entrenan cinco iteraciones que aíslan por
construcción tres palancas: adaptación textual, adaptación visual y
aumento de escala con imágenes externas genéricas. El objetivo
experimental, fijado a priori, es L3 \emph{accuracy}@1 $\geq 0{,}25$ en
\emph{zero-shot} sobre las $54$ clases L3 representadas en el test. La
arquitectura completa, los hiperparámetros, las particiones
\emph{case-aware} y los detalles de generación de \emph{captions},
costes y tiempos figuran en los Anexos~\ref{cap:anexoa}
(tabla~\ref{tab:spanderm-config}), \ref{cap:anexoc}
y~\ref{cap:anexoi}.

\subsection{Evolución v1$\to$v5}

La tabla~\ref{tab:spanderm-zs} reporta la \emph{accuracy}@1 sobre las
$101$ imágenes de test bajo régimen híbrido \emph{H6} (logits visuales
y textuales con pesos $w_v = 0{,}3$, $w_t = 0{,}7$ calibrados sobre
validación), uniforme a las cinco iteraciones para la comparativa
cross-checkpoint. Los niveles L1 y L2 \emph{zero-shot} sólo se reportan
para v4 y v5.

\begin{table}[H]
\centering
\footnotesize
\setlength{\tabcolsep}{3pt}
\renewcommand{\arraystretch}{0.95}
\caption{Clasificación \emph{zero-shot} de la rama contrastiva de
Dermapixel R0 sobre el
conjunto de test de $101$ imágenes (per-case estratificado por
L1) bajo régimen híbrido \emph{H6} homogéneo
($w_v = 0{,}3$, $w_t = 0{,}7$) aplicado a las cinco iteraciones.
La columna \emph{Objetivo} corresponde al umbral fijado en el
diseño experimental antes de observar resultados. Mejor cifra
por fila excluyendo la columna \emph{Objetivo} en negrita.
La trayectoria completa con configuración LoRA por iteración,
lift de recuperación y techo \emph{Linear Probing} figura en el
Anexo~\ref{cap:anexoi}
(tabla~\ref{tab:spanderm_iteraciones}).
Provenance: commit \texttt{d0190c3} del repositorio
\texttt{spanderm-clip}.}
\label{tab:spanderm-zs}
\begin{tabular}{lrrrrrrr}
\hline
Tarea                        & Random    & v1                  & v2                  & v3                  & v4                  & v5                  & Objetivo \\
\hline
L1 acc@1 ($4$ cls)            & $0{,}250$ & ---                 & ---                 & ---                 & $0{,}4257$          & $\mathbf{0{,}4554}$ & ---           \\
L2 acc@1 ($23$ cls)           & $0{,}043$ & ---                 & ---                 & ---                 & $0{,}1837$          & $\mathbf{0{,}2041}$ & ---           \\
\textbf{L3 acc@1 ($54$ cls)}  & $0{,}019$ & $0{,}1782$          & $0{,}1980$          & $0{,}1980$          & $\mathbf{0{,}2475}$ & $0{,}2079$          & $\geq 0{,}25$ \\
\hline
\end{tabular}
\end{table}

Las tres primeras iteraciones quedan empatadas en
$[0{,}1782,\,0{,}1980]$ ($\pm 2$ imágenes sobre $N = 101$). La cuarta
alcanza $0{,}2475$, a $0{,}0025$ del umbral objetivo ($0{,}25$, una
imagen de test, dentro del ruido). La quinta, ablación deliberada de
escala bajo idéntica LoRA que v4, desciende a $0{,}2079$ ($-4{,}0$\,pp).
El barrido de épocas que ilustra la estabilidad de v1$\to$v3 y la
síntesis tabular completa de las cinco iteraciones figuran en el
Anexo~\ref{cap:anexoi}
(figura~\ref{fig:spanderm-epoch-sweep},
tabla~\ref{tab:spanderm_iteraciones}).

\subsection{Hallazgos científicos}
\label{sec:spanderm-escala-especificidad}

El experimento arroja tres conclusiones diferenciadas.

\paragraph{(1) La adaptación textual aporta mejoras limitadas.}
Las iteraciones v1$\to$v3, que actúan sobre el codificador textual y
las \emph{captions} ricas, quedan empatadas en
$[0{,}1782,\,0{,}1980]$. El techo de \emph{Linear Probing} sobre las
\emph{features} $1\,024$-D del encoder visual congelado
($0{,}079$--$0{,}089$ en L3) ya es inferior al \emph{zero-shot} L3 de
v1/v2/v3 ($\sim 0{,}13$): ninguna intervención textual puede extraer
información que las \emph{features} visuales no contienen. El diagnóstico
mecanístico completo figura en el Anexo~\ref{cap:anexoi}
(tabla~\ref{tab:spanderm-lp-ceiling}).

\paragraph{(2) La adaptación visual es la principal fuente de mejora.}
El salto $+4{,}95$\,pp entre v3 y v4 bajo idéntico régimen \emph{H6}
($0{,}1980 \to 0{,}2475$) se atribuye a la ampliación de capacidad del
codificador visual (LoRA $r=32$ sobre doce bloques frente a $r=16$
sobre cuatro). El cuello de botella no reside en la formulación textual
ni en las \emph{captions} ricas, sino en la capacidad visual para
reesculpir las \emph{features} \emph{fine-grained} castellanas. Es la
palanca de mayor efecto del experimento.

\paragraph{(3) Ampliar con imágenes externas genéricas no mejora
(resultado negativo).} La quinta iteración, que añade $15\,128$
imágenes externas genéricas (PAD-UFES-20~\cite{padufes},
DermNet~\cite{dermnet}) con \emph{captions} sintéticas, reduce la
\emph{accuracy}@1 L3 a $0{,}2079$ ($-4{,}0$\,pp), resultado
\emph{flat-to-negative}. El encoder visual mejora aisladamente en
régimen visual puro (\emph{H1} $0{,}168 \to 0{,}238$, $+41{,}2\%$
relativo) y el techo \emph{Linear Probing} L2 sube $+54{,}4\%$ relativo,
pero el régimen híbrido L3 se degrada porque las \emph{captions}
externas no preservan el registro discursivo del archivo y erosionan la
alineación castellana sobre la que se calibra \emph{H6}. La cola larga
no se rellena: las $327$ clases L3 con $\leq 5$ imágenes permanecen sin
cobertura adicional (sólo ocho ganan cobertura). Superar el
\emph{zero-shot} L3 $\sim 0{,}25$ requiere anotación experta dirigida
\emph{in-domain} sobre las clases peor cubiertas, no escala genérica
(sección~\ref{sec:linea-active-learning}). El hallazgo es convergente
con la sensibilidad de DermLIP al tokenizador
(sección~\ref{sec:zs-results}): la alineación cross-modal de un dominio
clínico de alta especificidad lingüística es un recurso frágil.

Una ablación independiente con el modelo fundacional inglés DermLIP v2
($\approx 403\,000$ pares Derm1M~\cite{derm1m}) confirma que la palanca
limitante era la escala de la alineación cross-modal y no la
formulación textual: alcanza \emph{accuracy}@1 L3 $= 0{,}2871$ en
régimen \emph{H6} comparable ($+3{,}96$\,pp sobre el $0{,}2475$ de v4).
El análisis completo, con el eje idioma puro y el \emph{trade-off}
operativo (rendimiento bruto frente a privacidad, especificidad y
autonomía), figura en el Anexo~\ref{cap:anexoi}.

La diferencia entre regímenes es la esperada: la cabeza supervisada
cerrada con LoRA (sección~\ref{sec:dermapixel-spanderm-v0}) alcanza
Acc@1 $0{,}250$ sobre L2, frente a $0{,}133$ de la rama contrastiva
\emph{zero-shot open-class} (v2) sobre las $23$ clases efectivas del
test per-case (comparativa detallada en el Anexo~\ref{cap:anexoi},
tabla~\ref{tab:spanderm-l2-ceiling}). La ventaja operativa de la rama
contrastiva es aceptar consultas textuales libres en castellano sin
fijar el vocabulario, propiedad explotada en el módulo RAG del
prototipo (Anexo~\ref{cap:anexoh}).

\subsection{Recuperación multimodal}

El segundo eje es la recuperación cross-modal per-imagen sobre los
$107$ pares (imagen, caption rica) de validación, con script idéntico
entre versiones. La tabla~\ref{tab:spanderm-retrieval} reporta el
régimen \emph{rich-caption} (\texttt{llm\_case\_summary}, el más
representativo del lenguaje clínico real de la Dra.~R.~Taberner).

\begin{table}[H]
\centering
\footnotesize
\setlength{\tabcolsep}{4pt}
\renewcommand{\arraystretch}{0.95}
\caption{Recuperación cross-modal i2t R@$5$ de la rama contrastiva
de Dermapixel R0 sobre
los $107$ pares (imagen, caption rica) de validación, régimen
\emph{rich-caption} (\texttt{llm\_case\_summary}). La columna
\emph{Lift} reporta el factor multiplicativo sobre el baseline
aleatorio (\emph{pool} $= 107$, baseline $R@5 = 5/107 = 0{,}0467$).
v2 marca el pico histórico de alineación cross-modal castellano;
v5 muestra descenso coherente con la degradación documentada en
la tabla~\ref{tab:spanderm-zs} sobre el régimen híbrido. v1 y v3
no se reportan por insuficiencia de datos en el régimen
homogéneo cross-checkpoint. El \emph{lift} (factor de mejora sobre
el azar) se calcula como
$R@5_{\mathrm{abs}} / (k / N_{\mathrm{pool}})$, con
$k = 5$ y $N_{\mathrm{pool}} = 107$; normaliza respecto a la
dificultad intrínseca del \emph{pool} y constituye la métrica
honesta para comparar entre \emph{checkpoints}.}
\label{tab:spanderm-retrieval}
\begin{tabular}{lrrr}
\hline
Versión & i2t R@$5$ absoluto  & Referencia azar & \emph{Lift} sobre azar   \\
\hline
v2      & $\mathbf{0{,}654}$  & $0{,}047$       & $\mathbf{14{,}0\times}$  \\
v4      & $0{,}579$           & $0{,}047$       & $12{,}4\times$           \\
v5      & $0{,}505$           & $0{,}047$       & $10{,}8\times$           \\
\hline
\end{tabular}
\end{table}

La alineación cross-modal castellano-imagen alcanza su pico en v2:
i2t R@$5 = 0{,}654$ frente a la referencia aleatoria $0{,}047$ sobre
\emph{pool} de $107$, un \emph{lift} de $\mathbf{14{,}0\times}$
(figura~\ref{fig:spanderm-retrieval}), contribución empírica favorable
del experimento que sustenta el uso de la rama como módulo de
recuperación semántica. Las iteraciones posteriores conservan
\emph{lift} sustancial (v4 $12{,}4\times$, v5 $10{,}8\times$) pero
descienden respecto al pico, coherente con la degradación del régimen
híbrido (tabla~\ref{tab:spanderm-zs}): la LoRA visual añadida (v4) y las
\emph{captions} externas (v5) desplazan el espacio embebido respecto al
óptimo \emph{rich-caption} castellano de v2.

\begin{figure}[H]
\centering
\includegraphics[width=0.88\textwidth]{fig_spanderm_retrieval.png}
\caption{Recuperación cross-modal i2t R@$5$ de la rama contrastiva
de Dermapixel R0
sobre $107$ pares de validación, régimen \emph{rich-caption}.
Pico v2 de $\mathbf{14{,}0\times}$ de \emph{lift} sobre el azar
($0{,}047 \to 0{,}654$) atribuible al conjunto LoRA textual
$+$ \emph{captions} sintéticas ricas generadas mediante un modelo de
lenguaje grande. v4 y v5 conservan \emph{lift} sustancial
($12{,}4\times$ y $10{,}8\times$ respectivamente) pero descienden
respecto al pico, en coherencia con la degradación documentada
sobre el régimen híbrido.}
\label{fig:spanderm-retrieval}
\end{figure}

La calibración de los logits \emph{zero-shot} de la rama contrastiva
por \emph{temperature scaling} reduce el \emph{Expected Calibration
Error} (ECE) L3 un $\sim 97\%$ sobre v4 ($0{,}7262 \to 0{,}0213$,
$\tau = 0{,}4165$); el procedimiento, el diagrama de fiabilidad y la
limitación de convergencia sobre v5 figuran en el
Anexo~\ref{cap:anexoi}
(tabla~\ref{tab:spanderm-calibration},
figura~\ref{fig:spanderm-calibration}). El sistema v4 es candidato a
módulo de búsqueda textual castellano del prototipo DermApIxel
(Anexo~\ref{cap:anexoh}), con la calibración escalar $\tau = 0{,}4165$
como capa final.

\section{Interpretabilidad del SAE: alineamiento con conceptos
dermatoscópicos clásicos}
\label{sec:sae-derm7pt}

Esta sección analiza si las representaciones internas del modelo se
corresponden con los conceptos que usa un dermatólogo, base de la
interpretabilidad clínica. Las \emph{features} del autoencoder disperso
se alinean con los criterios dermatoscópicos clásicos, aunque forzar un
cuello de botella conceptual cuesta exactitud, como muestran los
experimentos sobre Derm7pt.

La línea de interpretabilidad mediante Sparse Autoencoders (SAE) del
Anexo~\ref{cap:anexof} introduce un SAE Large de $16\,384$
\emph{features} entrenado sobre las activaciones de PanDerm Large y
validado sobre los $48$ conceptos visuales de SkinCon~\cite{skincon}.
Esta sección extiende la validación a Derm7pt~\cite{kawahara2019}
($1\,011$ lesiones dermatoscópicas) y explota las anotaciones del
\emph{Seven-Point Checklist} (siete criterios por imagen) en tres
experimentos sobre DGX Spark GB10: E1 alineamiento \emph{single-feature}
entre las $16\,384$ activaciones y los siete criterios binarizados, E2
\emph{Concept Bottleneck Model} (CBM) jerárquico para predicción
melanoma/no-melanoma, y E3 cruce con los $16$ conceptos del Grupo A
propuestos a la Dra.~R.~Taberner. Los tres comparten Derm7pt completo
($N = 1\,011$, $24{,}5\%$ melanomas; \emph{splits} oficiales
$413$/$203$/$395$, $100$ melanomas en test). La activación del SAE
muestra $3\,214$ \emph{features} activas ($>5\%$), $2\,978$ inactivas
y $2\,717$ muy activas ($>50\%$), coherente con la sparsidad nominal.

E1 evalúa, para cada uno de los siete criterios binarizados, la AUROC
efectiva $\max(\mathrm{AUROC}, 1-\mathrm{AUROC})$ de cada una de las
$16\,384$ \emph{features} como predictor binario. La
tabla~\ref{tab:sae-d7pt-e1} reporta el AUROC efectivo de la mejor
\emph{feature} por criterio (\emph{top-1}) y el promedio de las diez
mejores (\emph{top-10}) sobre Derm7pt completo ($N = 1\,011$).

\begin{table}[H]
\centering
\footnotesize
\setlength{\tabcolsep}{4pt}
\renewcommand{\arraystretch}{0.95}
\caption{Alineamiento \emph{single-feature} entre las
$16\,384$ \emph{features} del SAE Large y los siete criterios
del \emph{Seven-Point Checklist} sobre Derm7pt
($N = 1\,011$). AUROC efectiva
$\max(\mathrm{AUROC}, 1-\mathrm{AUROC})$ de la mejor
\emph{feature} individual (\emph{top-1}) y promedio sobre las
diez mejores (\emph{top-10}). Mejor cifra agregada en negrita.}
\label{tab:sae-d7pt-e1}
\begin{tabular}{lrrrr}
\hline
Criterio                 & $N_{+}$ & $N_{-}$ & AUROC \emph{top-1} & AUROC \emph{top-10} \\
\hline
pigment\_network         & $611$ & $400$ & $0{,}754$ & $0{,}737$ \\
streaks                  & $358$ & $653$ & $0{,}688$ & $0{,}684$ \\
pigmentation             & $423$ & $588$ & $0{,}700$ & $0{,}657$ \\
regression\_structures   & $253$ & $758$ & $0{,}687$ & $0{,}673$ \\
dots\_and\_globules      & $782$ & $229$ & $0{,}681$ & $0{,}663$ \\
blue\_whitish\_veil      & $195$ & $816$ & $0{,}746$ & $0{,}721$ \\
vascular\_structures     & $188$ & $823$ & $\mathbf{0{,}823}$ & $\mathbf{0{,}782}$ \\
\hline
\end{tabular}
\end{table}

Las siete \emph{features} más alineadas individualmente alcanzan
AUROC efectivo en el intervalo $[0{,}681,\,0{,}823]$, con el
máximo sobre \emph{vascular structures} ($0{,}823$) y mínimo sobre
\emph{dots and globules} ($0{,}681$). La estabilidad entre el
\emph{top-1} y el \emph{top-10} (diferencia $\leq 0{,}043$ en
valor absoluto) indica que el alineamiento no se concentra en una
única \emph{feature} aislada sino que el SAE codifica los
conceptos mediante un pequeño subconjunto coherente de
activaciones.

E2 construye un CBM jerárquico de dos etapas: una cabeza logística
sobre las $16\,384$ activaciones estima la probabilidad de los siete
criterios, y un meta-clasificador logístico predice
melanoma/no-melanoma a partir de esas siete estimaciones. La
comparativa con el sondeo lineal directo sin cuello de botella figura
en la tabla~\ref{tab:sae-d7pt-cbm}.

\begin{table}[H]
\centering
\footnotesize
\setlength{\tabcolsep}{4pt}
\renewcommand{\arraystretch}{0.95}
\caption{Comparativa entre sondeo lineal directo sobre las
$16\,384$ \emph{features} del SAE Large y CBM jerárquico con
los siete criterios \emph{Seven-Point Checklist} como cuello de
botella. Conjunto de test Derm7pt
($N_{\mathrm{test}} = 395$, $100$ melanomas). AUROC en
negrita.}
\label{tab:sae-d7pt-cbm}
\begin{tabular}{lrrrrr}
\hline
Arquitectura                                & Acc@1     & BAcc      & AUROC              & W-F1      & Kappa     \\
\hline
Sondeo lineal directo (SAE $\to$ melanoma)  & $0{,}841$ & $0{,}735$ & $\mathbf{0{,}890}$ & $0{,}829$ & $0{,}527$ \\
CBM (SAE $\to$ $7$ conceptos $\to$ melanoma) & $0{,}825$ & $0{,}681$ & $\mathbf{0{,}832}$ & $0{,}801$ & $0{,}440$ \\
\hline
\end{tabular}
\end{table}

El paso de las activaciones crudas del SAE a través del cuello de
botella conceptual reduce la AUROC binaria melanoma de $0{,}890$ a
$0{,}832$ ($-5{,}8$\,pp), pérdida que documenta el compromiso entre
interpretabilidad y precisión. Los pesos del meta-clasificador
identifican \emph{blue whitish veil} y \emph{regression structures}
como los conceptos dominantes para la decisión melanoma, jerarquía
coherente con el peso clínico que el \emph{Seven-Point
Checklist}~\cite{kawahara2019} asigna a ambos criterios (criterios
mayores de $2$ puntos). El AUROC binario por concepto de la cabeza
intermedia y los pesos completos del meta-clasificador figuran en el
Anexo~\ref{cap:anexof} (tabla~\ref{tab:sae-d7pt-cbm-coefs}).

E3 cruza los $16$ conceptos dermatoscópicos del Grupo~A propuestos a
la Dra.~R.~Taberner (sección~\ref{sec:dermapixel}) con los siete
criterios del \emph{Seven-Point Checklist} sobre Derm7pt, reportando
para cada concepto con \emph{mapping} la AUROC efectiva de la mejor
\emph{feature} SAE (resultado completo en Anexo~\ref{cap:anexof},
tabla~\ref{tab:sae-d7pt-cross-rosa}). De los $16$ conceptos, $11$
admiten \emph{mapping} operativo y $7$ de ellos quedan cubiertos por al
menos una \emph{feature} SAE con AUROC efectivo $\geq 0{,}70$ (máximo
$0{,}926$ sobre vasos en horquilla, pese a sólo $15$ positivos sobre
$1\,011$, lo que ilustra que el SAE captura conceptos poco frecuentes
sin contaminación por la distribución mayoritaria). Los $5$ sin
\emph{mapping} (empedrado, paralelos surcos/crestas, lagunas
vasculares, vasos en corona) corresponden a estructuras no contempladas
en la escala original y definen el contenido de la anotación clínica
colaborativa.

\subsection*{Ajuste fino multitarea: cabezas conceptuales y
melanoma}
\label{sec:sae-d7pt-e4}

E4 replica el protocolo multitarea de
Kawahara et~al.~\cite{kawahara2019} sustituyendo la red convolucional
por PanDerm Large, con LoRA ($r = 16$) sobre los dos últimos bloques
(\emph{blocks.22-23}), idéntico al protocolo de Dermapixel R0
(sección~\ref{sec:dermapixel-spanderm-v0}). Sobre la representación
adaptada se conectan ocho cabezas paralelas (una por criterio del
\emph{Seven-Point Checklist} más una binaria de melanoma), optimizadas
con la suma equiponderada de entropías cruzadas
$\mathcal{L} = \sum_{i=1}^{7} \mathrm{CE}(\hat{c}_i, c_i) +
\mathrm{CE}(\hat{m}, m)$ y \emph{class weight} balanceado ($0{,}56$\,M
parámetros entrenables, $0{,}18\%$ del codificador). La
tabla~\ref{tab:sae-d7pt-e4-comp} compara E4 con el sondeo lineal
directo (E2 Direct LP) y el CBM jerárquico (E2 CBM) sobre el test
Derm7pt ($N_{\mathrm{test}} = 395$, $n_{\mathrm{mel}} \approx 98$),
con IC95\% por \emph{bootstrap}.

\begin{table}[H]
\centering
\footnotesize
\setlength{\tabcolsep}{4pt}
\renewcommand{\arraystretch}{0.95}
\caption{Comparativa sobre la tarea binaria melanoma/no-melanoma
entre sondeo lineal directo sobre las $16\,384$ \emph{features}
del SAE Large (E2 Direct LP), CBM jerárquico de siete conceptos
(E2 CBM) y \emph{fine-tuning} multitarea LoRA con ocho cabezas
paralelas (E4). Conjunto de test Derm7pt
($N_{\mathrm{test}} = 395$, $n_{\mathrm{mel}} \approx 98$), mismo
\emph{split} oficial. AUROC en negrita.}
\label{tab:sae-d7pt-e4-comp}
\begin{tabular}{lrrrrrr}
\hline
Método                                       & Params entr. & Acc@1     & BAcc      & AUROC              & W-F1      & Kappa     \\
\hline
E2 Direct LP (SAE $\to$ melanoma)             & $16$\,K       & $0{,}841$ & $0{,}735$ & $0{,}890$          & $0{,}829$ & $0{,}527$ \\
E2 CBM (SAE $\to$ $7$ conceptos $\to$ mel)    & $\sim 120$    & $0{,}825$ & $0{,}681$ & $0{,}832$          & $0{,}801$ & $0{,}440$ \\
\textbf{E4 FT multitarea LoRA}                & $0{,}56$\,M   & $\mathbf{0{,}841}$ & $\mathbf{0{,}804}$ & $\mathbf{0{,}891}$ & $\mathbf{0{,}843}$ & $\mathbf{0{,}590}$ \\
\hline
\end{tabular}
\end{table}

La cabeza melanoma del \emph{fine-tuning} multitarea alcanza AUROC
test $= 0{,}891$ (IC95\% $[0{,}856,\,0{,}927]$), que iguala al sondeo
lineal directo ($0{,}890$, $+0{,}1$\,pp, dentro del ruido del
\emph{bootstrap}) y mejora el CBM jerárquico en $+5{,}9$\,pp AUROC,
ofreciendo además predicción melanoma e inferencia sobre los siete
criterios en una sola pasada sin coste discriminativo apreciable. El
desglose por concepto (Anexo~\ref{cap:anexof},
tabla~\ref{tab:sae-d7pt-e4-concepts}) muestra que las cabezas
conceptuales del régimen multitarea pierden entre $1$ y $3$\,pp de
AUROC frente al sondeo lineal separado en seis de los siete conceptos
(máximo $-2{,}7$\,pp en \emph{regression structures}; dos mejoran
ligeramente): que la cabeza melanoma no sufra penalización medible
indica que el aprendizaje multitarea no degrada la tarea principal
cuando los conceptos auxiliares son taxonómicamente coherentes con
ella, como ocurre con el \emph{Seven-Point Checklist}, instrumento
diagnóstico de melanoma por construcción~\cite{kawahara2019}.

En síntesis, el SAE Large codifica conceptos dermatoscópicos clásicos
sin supervisión explícita y generaliza a Derm7pt sin reentrenamiento
(E1), el CBM jerárquico cuantifica el compromiso
interpretabilidad--precisión ($-5{,}8$\,pp AUROC, E2), el
\emph{fine-tuning} multitarea LoRA iguala al sondeo lineal directo en
melanoma preservando inferencia concepto-a-concepto (E4), y $7$ de los
$11$ conceptos del Grupo~A con \emph{mapping} quedan cubiertos por
\emph{features} SAE, definiendo los $5$ restantes el contenido de la
anotación clínica colaborativa (sección~\ref{sec:linea-derm7pt-sae}).
El conjunto refuerza la viabilidad del módulo M3 de interpretabilidad
del prototipo DermApIxel (Anexo~\ref{cap:anexoh}).

\section{Síntesis comparativa de métodos sobre DermapixelAI 1.0}
\label{sec:dermapixel-extra}
\label{sec:dermapixel-extra-sintesis}

Las secciones~\ref{sec:dermapixel-results}--\ref{sec:spanderm-clip}
evalúan codificadores de dominio (PanDerm Base/Large, DermLIP v2),
vision-lenguaje médicos (BiomedCLIP, SigLIP-SO400M) y las dos ramas de
Dermapixel R0 sobre DermapixelAI~1.0. El Anexo~\ref{cap:anexoj} añade
dos ablaciones complementarias ---codificadores generalistas no
médicos (DINOv2, CLIP-L) bajo sondeo lineal y dos arquitecturas de
inferencia (\emph{retrieval} $k$-NN con FAISS y \emph{zero-shot}
jerárquico en cascada)---, integradas en la síntesis. La
tabla~\ref{tab:dermapixel-extra-sintesis} consolida BAcc y AUROC por
nivel ontológico para el conjunto completo de métodos.

\begin{table}[H]
\centering
\footnotesize
\setlength{\tabcolsep}{4pt}
\renewcommand{\arraystretch}{0.95}
\caption{Síntesis de codificadores y arquitecturas evaluados sobre
DermapixelAI~1.0 (test fijo: $N = 36$ en L1/L2, $N = 28$ en L3).
La columna \emph{Ref.} remite a la sección de origen de las cifras
completas con sus intervalos de confianza. Mejor valor por columna
y nivel en negrita.}
\label{tab:dermapixel-extra-sintesis}
\begin{tabular}{llrrrrrr}
\hline
                                  &     & \multicolumn{2}{c}{L1} & \multicolumn{2}{c}{L2} & \multicolumn{2}{c}{L3} \\
Método                            & Ref.& BAcc & AUROC           & BAcc & AUROC           & BAcc & AUROC \\
\hline
PanDerm Base LP                   & §4.11 & $0{,}570$          & $0{,}697$
                                          & $0{,}221$          & $0{,}707$
                                          & $0{,}132$          & $0{,}735$ \\
PanDerm Large LP                  & §4.11 & $\mathbf{0{,}735}$ & $\mathbf{0{,}796}$
                                          & $0{,}265$          & $0{,}793$
                                          & $0{,}184$          & $0{,}813$ \\
DermLIP v2 LP                     & §4.11 & ---                & ---
                                          & ---                & $0{,}860$
                                          & ---                & --- \\
DINOv2 ViT-L LP                   & Anx.~\ref{cap:anexoj} & $0{,}537$          & $0{,}650$
                                          & $0{,}250$          & $0{,}718$
                                          & $0{,}211$          & $0{,}794$ \\
CLIP-L-336 LP                     & Anx.~\ref{cap:anexoj} & $0{,}556$          & $0{,}658$
                                          & $0{,}279$          & $\mathbf{0{,}826}$
                                          & $0{,}184$          & $\mathbf{0{,}827}$ \\
Dermapixel R0 LoRA                & §4.14 & ---                & ---
                                          & $\mathbf{0{,}363}$ & $0{,}855$
                                          & ---                & --- \\
FT parcial denso L1$+$L2$+$L3      & §4.11 & $0{,}622$          & $0{,}715$
                                          & $0{,}324$          & $0{,}852$
                                          & $0{,}184$          & $0{,}804$ \\
$k$-NN FAISS PanDerm Large         & Anx.~\ref{cap:anexoj} & $0{,}686$          & $0{,}748$
                                          & $0{,}265$          & $0{,}700$
                                          & $\mathbf{0{,}210}$ & $0{,}719$ \\
\hline
\end{tabular}
\end{table}

El mejor método sobre DermapixelAI~1.0 depende del nivel
ontológico y de la métrica. Las cifras de esta síntesis se reportan
sobre el test fijo ($N = 36$ en L1/L2, $N = 28$ en L3), por lo que las
jerarquías entre métodos en L2 y L3 deben leerse como indicio y
remitirse a la validación cruzada agrupada por caso
(tabla~\ref{tab:dermapixel-cv}) como estimación defendible del
rendimiento esperado del sondeo lineal congelado. Sobre L1, el sondeo
lineal sobre PanDerm Large encabeza simultáneamente BAcc ($0{,}735$) y
AUROC ($0{,}796$); sobre L2, Dermapixel R0 LoRA encabeza la BAcc
($0{,}363$) y CLIP-L-336 LP encabeza la AUROC ($0{,}826$); sobre
L3, el \emph{retrieval} $k$-NN FAISS sobre embeddings de PanDerm
Large encabeza la BAcc ($0{,}210$) y CLIP-L-336 LP encabeza la
AUROC ($0{,}827$). La discriminación binaria melanoma/no-melanoma
sobre Derm7pt (sección~\ref{sec:sae-derm7pt}, fuera de
DermapixelAI~1.0) alcanza AUROC $0{,}891$ con el \emph{fine-tuning}
multitarea LoRA E4 sobre las activaciones del SAE. La ausencia de
un método único óptimo refuerza la lectura interpretativa de la
sección~\ref{sec:disc-especializacion}: la elección del
codificador y del clasificador debe condicionarse al nivel
ontológico objetivo y al criterio operativo predominante
(discriminación robusta, ranking probabilístico o recuperación
para soporte explicativo).


\section{Aplicación de seguridad: ensemble para detección de melanoma}
\label{sec:ensemble-results}

Más allá de la comparación de modelos individuales, resulta
relevante explorar si la complementariedad entre arquitecturas puede
explotarse para maximizar la sensibilidad diagnóstica. Esta cuestión
es especialmente importante en melanoma, donde el coste clínico de
un falso negativo es potencialmente muy superior al de un falso
positivo. Por ello se evalúa un \emph{ensemble} orientado
explícitamente a maximizar el \emph{recall} de melanoma sobre
HAM10000.

Se combinan tres clasificadores independientes sobre el
\emph{split} de test de HAM10000 ($N = 1\,232$ imágenes,
$70$ melanomas):
M1 (PanDerm Large \emph{fine-tuned} sobre HAM10000), SigLIP-Large
SO400M con sondeo lineal sobre los \emph{embeddings} de
$1\,152$ dimensiones, y M7 (clasificador unificado sobre el
\emph{merged43} del dataset armonizado, descrito en el
Anexo~\ref{cap:anexoh}).

\begin{table}[H]
\centering
\footnotesize
\setlength{\tabcolsep}{4pt}
\renewcommand{\arraystretch}{0.95}
\caption{Ensemble sobre HAM10000 ($N = 1\,232$, $70$ melanomas).
\emph{Recall} \emph{top-3}: la clase melanoma aparece entre las
tres primeras del ranking. FP\,mel: falsos positivos melanoma
(imágenes no-melanoma marcadas).}
\label{tab:ensemble}
\begin{tabular}{lrrrrrr}
\hline
Configuración                   & Acc           & BAcc          & AUROC         & Mel rec.\ top-1 & Mel rec.\ top-3 & FP\,mel \\
\hline
M1 (PanDerm FT)                 & $0{,}919$     & $0{,}813$     & $0{,}986$     & $0{,}714$ (50/70) & $0{,}971$ (68/70) & $36$ \\
SigLIP-L LP                     & $0{,}900$     & $0{,}776$     & $0{,}971$     & $0{,}614$ (43/70) & $0{,}986$ (69/70) & $31$ \\
Avg M1 + SigLIP                 & $0{,}917$     & $0{,}836$     & $0{,}986$     & $0{,}686$         & $0{,}986$         & $32$ \\
MaxMel M1 + SigLIP              & $0{,}909$     & $0{,}832$     & $0{,}985$     & $\mathbf{0{,}771}$ & $\mathbf{1{,}000}$ & $50$ \\
OR top-3 M1 + SigLIP            & ---           & ---           & ---           & ---               & $\mathbf{1{,}000}$ (70/70) & $885$ \\
\hline
\end{tabular}
\end{table}

La combinación OR \emph{top-3} de M1 y SigLIP alcanza un
\emph{recall} de melanoma del $100\%$ ($70$ de $70$), pero a costa
de $885$ falsos positivos sobre las $1\,232$ imágenes del test: la
sensibilidad máxima se obtiene marcando como sospechosa cerca del
$72\%$ del conjunto, lo que reduce drásticamente la precisión y
acota su utilidad a un escenario de \emph{cribado}, nunca de
decisión diagnóstica.

El principal hallazgo de esta sección, sin embargo, no es ese
$100\%$ ---una cifra trivialmente alcanzable ampliando el margen de
sospecha de cualquier clasificador--- sino la evidencia de
complementariedad entre una arquitectura especializada (M1, PanDerm
\emph{fine-tuned} \emph{in-domain}) y una generalista (SigLIP con
sondeo lineal): ambas recuperan melanomas distintos. Los dos
melanomas que M1 omite en \emph{top-3} (ISIC\_0024886.jpg e
ISIC\_0030360.jpg) sí los recupera SigLIP, y el único melanoma que
SigLIP no sitúa en \emph{top-3} sí lo recupera M1. Este resultado
indica que los errores de ambos modelos no están completamente
correlacionados ---condición necesaria para que una estrategia de
\emph{ensemble} aporte un beneficio clínico real, y no una mera
redundancia--- y sugiere que cada arquitectura explota
representaciones parcialmente diferentes de la lesión. M7, evaluado
sobre el subconjunto de $198$ imágenes con predicción simultánea de
los tres modelos ($12$ melanomas), no recupera melanomas adicionales
sobre el par, pero aporta cobertura sobre las $43$ clases
ontológicas del dataset armonizado.

Tres caveats acompañan la interpretación de este resultado: (i)~el
\emph{recall} del $100\%$ es \emph{top-3}, no \emph{top-1};
(ii)~el tamaño muestral de melanomas es $N = 70$, con intervalos
de confianza al $95\%$ amplios; (iii)~no existe validación
prospectiva sobre cohorte hospitalaria. El resultado se interpreta
como prueba de concepto para un escenario de cribado oportunista,
no como sustituto del juicio clínico.
